\documentclass[11pt]{article}
\usepackage[margin=1in]{geometry}
\usepackage{amsmath}
\usepackage{amssymb}
\usepackage{array}
\usepackage{booktabs}
\usepackage{caption}
\usepackage{float}
\usepackage{graphicx}
\usepackage{longtable}
\usepackage{microtype}
\usepackage{pdflscape}
\usepackage{subcaption}
\usepackage{threeparttable}
\usepackage{hyperref}
\usepackage{xurl}
\usepackage{xcolor}
\usepackage{placeins}

\hypersetup{colorlinks=true,linkcolor=blue,urlcolor=blue,citecolor=blue}
\captionsetup{font=small,labelfont=bf}
\subcaptionsetup{font=small,labelfont=bf,justification=centering}
\setlength{\parindent}{1.2em}
\setlength{\parskip}{0.15em}
\setlength{\emergencystretch}{2em}

\newcommand{\figdir}{figuresNtables}
\newcommand{\panelwidth}{0.76\textwidth}

\title{AI Exposure and Labor-Market Outcomes in Spain\\
\large Estimates Results Report, Version 1}
\author{}
\date{\today}

\begin{document}
\maketitle

\setcounter{section}{3}
\section{Results}
\label{sec:v1_results}

\subsection{Baseline Results}
\label{sec:v1_baseline}

Figure~\ref{fig:v1_descriptive} describes the evolution of the two labor-market outcomes before turning to the regression evidence. Registered unemployment falls in all three exposure groups, but the decline becomes slower among highly exposed occupations after November 2022. The separation is gradual rather than immediate, as one would expect if firms adopted generative AI over time. New contract registrations are much more volatile and show no stable ordering by exposure. The similar pre-release paths are reassuring, but they are not sufficient for identification because exposure groups may differ in their recovery from the pandemic and in other occupation-specific shocks.

\begin{figure}[H]
    \centering
    \caption{Descriptive patterns by AI exposure}
    \label{fig:v1_descriptive}
    \begin{subfigure}{\panelwidth}
        \centering
        \includegraphics[width=\textwidth]{\figdir/Figure1_panelA.png}
        \caption{Number of registered unemployed}
    \end{subfigure}

    \vspace{0.30cm}

    \begin{subfigure}{\panelwidth}
        \centering
        \includegraphics[width=\textwidth]{\figdir/Figure1_panelB.png}
        \caption{Number of new contracts}
    \end{subfigure}

    \begin{minipage}{0.92\textwidth}
    \footnotesize \emph{Notes:} The figure plots equal-weighted means of
    occupation-specific outcome indices in the total CNO4 panel. Each
    occupation's outcome is normalized to November 2022 = 100 before averaging
    within exposure group. Zero-exposure occupations have a nearest-neighbor
    score of zero; middle-exposure occupations have a positive score below
    0.1169; and high-exposure occupations have a score at or above 0.1169, the
    75th percentile of exposure. The vertical dashed line marks November 2022,
    the public release of ChatGPT. Registered unemployment is a stock; newly
    registered contracts are monthly flows. Panel B excludes two occupations
    with zero contracts in the reference month. The figure is descriptive and
    does not report causal estimates.
    \end{minipage}
\end{figure}

The preferred specification asks whether outcomes diverged after ChatGPT's release across occupations with different exposure, while comparing occupations within the same broad CNO1 family in every month. This restriction is economically important: it removes shocks shared by related occupations, such as changes in sectoral demand or the pace of the post-pandemic recovery, without requiring those shocks to follow a parametric trend.\footnote{For CNO4 occupation \(j\), CNO1 family \(g(j)\), and event time \(k\), the event-study regression is \(Y_{jt}=\alpha_j+\lambda_{g(j),t}+\sum_{k\neq -1}\beta_k(A_j/0.10)\mathbf{1}\{t-T_0=k\}+\varepsilon_{jt}\). The event window runs from \(-21\) to \(40\), November 2022 is \(T_0\), and October 2022 is omitted. Occupation fixed effects absorb permanent differences across CNO4 occupations, CNO1-by-month fixed effects absorb shocks common to occupations in the same broad family, and standard errors are clustered by CNO4.}

Figure~\ref{fig:v1_eventstudy} reports the preferred dynamic estimates. In panel (a), the unemployment gradient is close to zero around the release date and becomes positive only after several months. A 10 percentage-point increase in exposure is associated with roughly 1--2 percent more registered unemployment over much of the post-treatment period. The estimates remain positive near the end of the sample, but their confidence intervals widen and frequently include zero. Panel (b) shows no comparable persistence in new contracts: the coefficients fluctuate around zero and are imprecisely estimated. Taken together, the panels are consistent with a gradual increase in the stock of registered unemployment rather than a discrete contraction in the monthly flow of hiring. They do not, by themselves, distinguish displacement from slower exits out of unemployment.

The pre-treatment evidence warrants caution. For registered unemployment, the joint-nullity and flatness tests reject over event times \(-21\) through \(-2\) (both \(p=0.003\)), but not over the more recent window \(-10\) through \(-2\) (\(p=0.367\) and \(0.280\)). The early window accounts for the full-period rejection. For new contracts, the corresponding tests reject over the full window (\(p=0.032\) and \(0.045\)) and are borderline over the recent window (\(p=0.049\) and \(0.035\)). Allowing errors to be correlated within CNO3 groups does not overturn the broad conclusion. Appendix Table~\ref{tab:v1_pretrends} reports all windows and clustering choices. The CNO1-by-month effects improve the economic comparability of occupations, especially in the period closest to treatment, but they do not eliminate the possibility that differential recovery paths contaminate the post-treatment gradient. We therefore interpret the estimates as disciplined but still assumption-dependent evidence.

\begin{figure}[H]
    \centering
    \caption{Impact of AI exposure on labor market outcomes. Event-study estimates.}
    \label{fig:v1_eventstudy}
    \begin{subfigure}{\panelwidth}
        \centering
        \includegraphics[width=\textwidth]{\figdir/Figure2_panelA.png}
        \caption{Number of registered unemployed}
    \end{subfigure}

    \vspace{0.30cm}

    \begin{subfigure}{\panelwidth}
        \centering
        \includegraphics[width=\textwidth]{\figdir/Figure2_panelB.png}
        \caption{Number of new contracts}
    \end{subfigure}

    \begin{minipage}{0.92\textwidth}
    \footnotesize \emph{Notes:} The figure reports estimates of the preferred
    event-study specification for the total CNO4 occupation panel. The
    dependent variables are the logarithm of registered unemployed workers in
    panel (a) and the logarithm of new contracts in panel (b). Each marker is
    the marginal effect of a 10 percentage-point increase in AI exposure at the
    indicated event month. October 2022 (event time \(-1\)) is the omitted
    reference month; the vertical dashed line marks November 2022, the public
    release of ChatGPT. All regressions include CNO4 and CNO1-by-year-month
    fixed effects. Shaded areas report 95 percent confidence intervals based on
    standard errors clustered by CNO4.
    \end{minipage}
\end{figure}

\setcounter{table}{1}
\begin{table}[!htbp]
\centering
\caption{Impact of AI exposure on labor market outcomes}
\label{tab:v1_main_effects}
\begin{threeparttable}
\scriptsize
\setlength{\tabcolsep}{3.5pt}
\begin{tabular}{lcccccc}
\toprule
& \multicolumn{3}{c}{\# of registered unemployed} & \multicolumn{3}{c}{\# of new contracts} \\
\cmidrule(lr){2-4}\cmidrule(lr){5-7}
& (1) & (2) & (3) & (4) & (5) & (6) \\
\midrule
AI exposure $\times$ adjustment period & 0.019$^{***}$ & 0.012$^{**}$ & 0.012$^{*}$ & 0.033$^{***}$ & 0.014 & 0.014 \\
 & (0.005) & (0.006) & (0.007) & (0.008) & (0.009) & (0.009) \\
Impact of a 10 pp increase (percent) & 1.9 & 1.2 & 1.2 & 3.3 & 1.4 & 1.4 \\
\addlinespace
AI exposure $\times$ later period & 0.038$^{***}$ & 0.017 & 0.017 & 0.029$^{***}$ & 0.013 & 0.013 \\
 & (0.010) & (0.012) & (0.013) & (0.010) & (0.011) & (0.012) \\
Impact of a 10 pp increase (percent) & 3.8 & 1.7 & 1.7 & 2.9 & 1.3 & 1.3 \\
$p$-value: equal phase effects & 0.001 & 0.486 & 0.529 & 0.488 & 0.807 & 0.841 \\
\midrule
CNO1 $\times$ year-month FE & No & Yes & Yes & No & Yes & Yes \\
Clustered standard errors & CNO4 & CNO4 & CNO3 & CNO4 & CNO4 & CNO3 \\
Observations & 31,122 & 31,122 & 31,122 & 30,787 & 30,786 & 30,786 \\
\bottomrule
\end{tabular}
\begin{tablenotes}[flushleft]
\footnotesize
\item \emph{Notes:} Entries are marginal effects from occupation-month regressions. The adjustment period covers November 2022 through November 2024 (event times 0--24); the later period covers December 2024 through March 2026 (event times 25--40). All pre-treatment months form the omitted category. The dependent variables are logarithms. Columns 1 and 4 include CNO4 and year-month fixed effects; columns 2, 3, 5, and 6 include CNO4 and CNO1-by-year-month fixed effects. Exposure is divided by 0.10, so coefficients correspond to a 10 percentage-point increase. Standard errors are clustered as indicated. The impact rows report 100 times the coefficient and carry no significance symbols. The equality row reports the p-value for the null that the two phase effects are equal. These are marginal effects, not ATTs. $^{***}p<0.01$, $^{**}p<0.05$, and $^{*}p<0.10$.
\end{tablenotes}
\end{threeparttable}
\end{table}

\FloatBarrier

Table~\ref{tab:v1_main_effects} summarizes the post-treatment path in two intervals. The adjustment period covers November 2022 through November 2024, when firms were learning about and beginning to deploy generative AI. The later period covers December 2024 through March 2026, when adoption had more time to affect staffing and job search.\footnote{The table estimates \(Y_{jt}=\alpha_j+\lambda_{g(j),t}+\beta_A(A_j/0.10)\mathbf{1}\{0\leq e_t\leq24\}+\beta_L(A_j/0.10)\mathbf{1}\{25\leq e_t\leq40\}+\varepsilon_{jt}\), where \(\lambda_{g(j),t}\) denotes common month effects or CNO1-by-month effects. All pre-treatment months form the omitted category. The table coefficients compare each interval with the average pre-treatment exposure gradient; Figure~\ref{fig:v1_eventstudy} instead normalizes each monthly coefficient to October 2022.}

The benchmark in column 1 implies 1.9 percent more registered unemployment during the adjustment period and 3.8 percent more in the later period for a 10 percentage-point increase in exposure. Restricting identification to within-CNO1 comparisons attenuates the estimates to 1.2 and 1.7 percent in column 2. The adjustment-period coefficient is statistically significant at the 5 percent level; the later coefficient is positive but less precisely estimated. The hypothesis of equal effects across the two intervals cannot be rejected (\(p=0.486\)). Clustering at CNO3 in column 3 leaves the point estimates unchanged and modestly increases the standard errors. The attenuation after adding CNO1-by-month effects is economically informative: part of the unconditional relationship reflects shocks shared by broad occupational families rather than variation between comparable CNO4 occupations.

The contract estimates display a similar attenuation. The benchmark coefficients in column 4 are 3.3 and 2.9 percent, while the preferred estimates in column 5 are 1.4 and 1.3 percent. Neither preferred estimate is statistically distinguishable from zero, and the two intervals are nearly identical (\(p=0.807\)). The baseline evidence therefore points to a modest increase in registered unemployment in more exposed occupations, but the effect is estimated imprecisely and is not accompanied by a persistent aggregate response in new contracts.

The absence of a CNO4 labor-force denominator prevents a direct occupational unemployment rate from being constructed. The estimates should therefore be interpreted as effects on registered unemployment outcomes rather than as effects on unemployment rates. This limitation does not invalidate the within-CNO4 fixed-effects design, but it prevents the results from being expressed as changes in the share of the labor force that is unemployed.

\subsubsection{Robustness Checks}
\label{sec:v1_robustness}

Appendix Table~\ref{tab:v1_robustness} examines whether the baseline result is an artifact of outcome coding or exposure construction. Replacing \(\log(Y)\) with \(\log(Y+1)\) leaves the preferred unemployment estimates at 1.2 percent in the adjustment period and 1.7 percent later. Cosine-weighted exposure produces somewhat larger estimates of 1.8 and 2.9 percent. Appendix Table~\ref{tab:v1_jev_robustness} applies the same long-difference specifications to three Jev-based exposure measures. Appendix Figure~\ref{fig:v1_exposure_correlations} compares the occupation rankings across measures. Appendix Table~\ref{tab:v1_robustness_pretrends} and Figures~\ref{fig:v1_robust_u}--\ref{fig:v1_jev_robust_c} report the corresponding event-study diagnostics.

The binary design in Appendix Table~\ref{tab:v1_binary_treatment} classifies occupations above the 75th-percentile cutoff, 0.1169, as treated and retains all remaining occupations as controls. The preferred specification yields 3.9 percent more registered unemployment during the adjustment period and 5.8 percent more later, although the later estimate is imprecise and the two interval effects are not statistically different. Contract estimates remain close to zero. The larger binary coefficients are consistent with concentration of the response in the upper exposure tail, but discretization discards dose information and changes the estimand. Appendix Table~\ref{tab:v1_binary_pretrends} and Figure~\ref{fig:v1_binary_event} report the corresponding diagnostics.

Appendix Table~\ref{tab:v1_province} disaggregates the panel by province and absorbs province-specific monthly shocks. The preferred unemployment estimates are 0.9 percent during the adjustment period and 1.0 percent later; only the former is marginally distinguishable from zero. Contract estimates are approximately 0.9 percent in both intervals and are imprecise. The smaller coefficients suggest that regional composition contributes to the aggregate gradient, although the direction remains positive. Appendix Figure~\ref{fig:v1_province_event} reports the monthly estimates.

No single CNO1 family drives the result: the leave-one-family-out estimates in Appendix Figure~\ref{fig:v1_leave_one_out} remain positive and close to the full-sample estimate. Appendix Table~\ref{tab:v1_trailing_contracts} and Figure~\ref{fig:v1_trailing_contracts} replace monthly contracts with registrations accumulated over the preceding 12 months. The preferred estimates, 1.5 percent during adjustment and 0.9 percent later, are not statistically different from zero. This measure smooths hiring flows but does not create a stock of active contracts, so it is a robustness check rather than a substitute outcome.

\subsection{Heterogeneity}
\label{sec:v1_heterogeneity}

\begin{table}[p]
\centering
\caption{Heterogeneity in the effects of AI exposure}
\label{tab:v1_heterogeneity}
\begin{threeparttable}
\scriptsize
\renewcommand{\arraystretch}{0.84}
\setlength{\tabcolsep}{3.0pt}
\begin{tabular}{@{}lcccccc@{}}
\toprule
& \multicolumn{3}{c}{\# of registered unemployed} & \multicolumn{3}{c}{\# of new contracts} \\
\cmidrule(lr){2-4}\cmidrule(lr){5-7}
& (1) & (2) & (3) & (4) & (5) & (6) \\
\midrule
\multicolumn{7}{l}{\textbf{Panel A. Age: under 30}} \\
AI exposure $\times$ adjustment period & 0.008 & 0.015 & 0.015 & 0.033$^{***}$ & 0.022$^{**}$ & 0.022$^{**}$ \\
 & (0.009) & (0.010) & (0.012) & (0.009) & (0.010) & (0.010) \\
Impact of a 10 pp increase (percent) & 0.8 & 1.5 & 1.5 & 3.3 & 2.2 & 2.2 \\
\addlinespace
AI exposure $\times$ later period & 0.028$^{**}$ & 0.022 & 0.022 & 0.027$^{**}$ & 0.026$^{**}$ & 0.026$^{*}$ \\
 & (0.014) & (0.016) & (0.019) & (0.011) & (0.013) & (0.015) \\
Impact of a 10 pp increase (percent) & 2.8 & 2.2 & 2.2 & 2.7 & 2.6 & 2.6 \\
$p$-value: equal phase effects & 0.005 & 0.461 & 0.522 & 0.290 & 0.519 & 0.591 \\
Observations & 30,493 & 30,493 & 30,493 & 30,160 & 30,139 & 30,139 \\
\addlinespace[1pt]
\multicolumn{7}{l}{\textbf{Panel B. Age: 30--39}} \\
AI exposure $\times$ adjustment period & 0.037$^{***}$ & 0.017$^{*}$ & 0.017 & 0.030$^{***}$ & 0.013 & 0.013 \\
 & (0.008) & (0.009) & (0.010) & (0.008) & (0.009) & (0.009) \\
Impact of a 10 pp increase (percent) & 3.7 & 1.7 & 1.7 & 3.0 & 1.3 & 1.3 \\
\addlinespace
AI exposure $\times$ later period & 0.075$^{***}$ & 0.034$^{**}$ & 0.034$^{**}$ & 0.031$^{***}$ & 0.017 & 0.017 \\
 & (0.014) & (0.015) & (0.017) & (0.010) & (0.011) & (0.013) \\
Impact of a 10 pp increase (percent) & 7.5 & 3.4 & 3.4 & 3.1 & 1.7 & 1.7 \\
$p$-value: equal phase effects & 0.000 & 0.023 & 0.040 & 0.809 & 0.439 & 0.531 \\
Observations & 30,776 & 30,776 & 30,776 & 30,270 & 30,255 & 30,255 \\
\addlinespace[1pt]
\multicolumn{7}{l}{\textbf{Panel C. Age: 40 or older}} \\
AI exposure $\times$ adjustment period & 0.017$^{***}$ & 0.009 & 0.009 & 0.028$^{***}$ & 0.003 & 0.003 \\
 & (0.004) & (0.005) & (0.006) & (0.009) & (0.009) & (0.010) \\
Impact of a 10 pp increase (percent) & 1.7 & 0.9 & 0.9 & 2.8 & 0.3 & 0.3 \\
\addlinespace
AI exposure $\times$ later period & 0.032$^{***}$ & 0.012 & 0.012 & 0.026$^{**}$ & -0.001 & -0.001 \\
 & (0.009) & (0.011) & (0.013) & (0.010) & (0.011) & (0.013) \\
Impact of a 10 pp increase (percent) & 3.2 & 1.2 & 1.2 & 2.6 & -0.1 & -0.1 \\
$p$-value: equal phase effects & 0.006 & 0.604 & 0.644 & 0.615 & 0.538 & 0.620 \\
Observations & 31,099 & 31,099 & 31,099 & 30,442 & 30,435 & 30,435 \\
\addlinespace[1pt]
\multicolumn{7}{l}{\textbf{Panel D. Gender: men}} \\
AI exposure $\times$ adjustment period & 0.017$^{***}$ & 0.008 & 0.008 & 0.032$^{***}$ & 0.014 & 0.014 \\
 & (0.005) & (0.007) & (0.008) & (0.009) & (0.009) & (0.010) \\
Impact of a 10 pp increase (percent) & 1.7 & 0.8 & 0.8 & 3.2 & 1.4 & 1.4 \\
\addlinespace
AI exposure $\times$ later period & 0.035$^{***}$ & 0.016 & 0.016 & 0.032$^{***}$ & 0.016 & 0.016 \\
 & (0.010) & (0.012) & (0.014) & (0.011) & (0.012) & (0.013) \\
Impact of a 10 pp increase (percent) & 3.5 & 1.6 & 1.6 & 3.2 & 1.6 & 1.6 \\
$p$-value: equal phase effects & 0.003 & 0.312 & 0.364 & 0.977 & 0.819 & 0.840 \\
Observations & 31,058 & 31,058 & 31,058 & 30,535 & 30,530 & 30,530 \\
\addlinespace[1pt]
\multicolumn{7}{l}{\textbf{Panel E. Gender: women}} \\
AI exposure $\times$ adjustment period & 0.014$^{**}$ & 0.010 & 0.010 & 0.022$^{**}$ & 0.011 & 0.011 \\
 & (0.006) & (0.007) & (0.008) & (0.009) & (0.009) & (0.009) \\
Impact of a 10 pp increase (percent) & 1.4 & 1.0 & 1.0 & 2.2 & 1.1 & 1.1 \\
\addlinespace
AI exposure $\times$ later period & 0.030$^{***}$ & 0.010 & 0.010 & 0.010 & 0.009 & 0.009 \\
 & (0.010) & (0.012) & (0.014) & (0.011) & (0.012) & (0.013) \\
Impact of a 10 pp increase (percent) & 3.0 & 1.0 & 1.0 & 1.0 & 0.9 & 0.9 \\
$p$-value: equal phase effects & 0.010 & 0.995 & 0.995 & 0.024 & 0.786 & 0.825 \\
Observations & 30,813 & 30,813 & 30,813 & 30,207 & 30,188 & 30,188 \\
\midrule
CNO1 $\times$ year-month FE & No & Yes & Yes & No & Yes & Yes \\
Clustered standard errors & CNO4 & CNO4 & CNO3 & CNO4 & CNO4 & CNO3 \\
\bottomrule
\end{tabular}
\begin{tablenotes}[flushleft]
\scriptsize
\item \emph{Notes:} Each panel reports subgroup-specific marginal effects for the adjustment period (event times 0--24) and later period (event times 25--40), relative to all pre-treatment months. Age categories are under 30, 30--39, and 40 or older. Columns 1 and 4 include occupation and year-month fixed effects; the remaining columns replace common month effects with CNO1-by-year-month effects. Standard errors are clustered as indicated. Exposure is measured in 10 percentage-point units. Equality rows test whether the two phase effects are equal. $^{***}p<0.01$, $^{**}p<0.05$, and $^{*}p<0.10$.
\end{tablenotes}
\end{threeparttable}
\end{table}


Table~\ref{tab:v1_heterogeneity} shows that the average response masks economically relevant differences by age. In the preferred within-CNO1 specification, a 10 percentage-point increase in exposure is associated with 1.5 percent more registered unemployment for workers under 30, 1.7 percent for ages 30--39, and 0.9 percent for workers age 40 or older during the adjustment period. The later-period estimates are 2.2, 3.4, and 1.2 percent. The estimate for ages 30--39 is individually significant in the later period, but separate significance levels do not establish that effects differ across groups. The pooled tests in Appendix Table~\ref{tab:v1_age_cross_group_tests} do not reject equality across all three groups (\(p=0.105\)); they do reject equality between ages 30--39 and age 40 or older (\(p=0.035\)), but not between ages 30--39 and workers under 30 (\(p=0.320\)). The age concentration is therefore suggestive rather than definitive.

The contract results provide a complementary view of adjustment. For workers under 30, the preferred coefficients are 2.2 percent during adjustment and 2.6 percent later, both statistically significant. The corresponding estimates for ages 30--39 are 1.3 and 1.7 percent, while those for workers age 40 or older are close to zero. The pooled tests reject equality across age groups in both periods. Younger workers may therefore experience more labor-market turnover, while the 30--39 group exhibits the clearest later accumulation in registered unemployment. The data do not observe worker transitions, however, so this interpretation cannot distinguish job creation, repeated short contracts, or slower exits from unemployment.

The age results should be interpreted as heterogeneity in occupation-by-age cells, rather than as evidence that individual workers in one age group are causally displaced. Early-to-mid-career workers may occupy tasks and career positions that are more sensitive to organizational adoption, but the present data cannot distinguish this mechanism from differences in job mobility, hiring, or unemployment duration. The pooled tests and subgroup pre-trend diagnostics therefore make the age pattern suggestive rather than definitive. Appendix Table~\ref{tab:v1_heterogeneity_pretrends} and Figures~\ref{fig:v1_heterogeneity_u}--\ref{fig:v1_heterogeneity_c} report the subgroup diagnostics and dynamic paths.\footnote{SEPE omits its CNO4-by-age breakdown for May 2024. We recover uniquely identified integer cells by inverting each June level and the corresponding June-over-May percentage change, with the observed May CNO4 total used as an adding-up validation. Cells that are not uniquely identified remain missing.}

\subsection{Heterogeneity by Occupational Feminization}
\label{sec:v1_feminization}

The occupational-feminization analysis asks whether exposure has different consequences in occupations with historically different gender composition. The predetermined index is the female share of EPA employment in the corresponding CNO2 group, averaged over 2017--2019. We split occupations at its median, 0.303, and allow the exposure gradient to differ on either side of that threshold within the preferred specification. This is an occupation-level task-composition exercise; it does not estimate different effects for individual men and women.\footnote{The estimating equation is \(\ln Y_{jt}=\alpha_j+\lambda_{g(j),t}+\sum_{p\in\{A,L\}}\sum_{h\in\{B,H\}}\beta_{ph}(A_j/0.10)\mathbf{1}\{e_t\in p\}\mathbf{1}\{F_j\in h\}+\varepsilon_{jt}\), where \(A\) and \(L\) denote the adjustment and later periods, \(F_j\) is the predetermined feminization index, and the pre-treatment period is omitted. Standard errors are clustered by CNO4.}

\begin{table}[!htbp]
\centering
\caption{AI exposure effects by occupational feminization}
\label{tab:v1_feminization_median_split}
\begin{threeparttable}
\small
\setlength{\tabcolsep}{5pt}
\begin{tabular}{@{}p{0.58\textwidth}>{\centering\arraybackslash}p{0.17\textwidth}>{\centering\arraybackslash}p{0.17\textwidth}@{}}
\toprule
& (1) & (2) \\
& \# of registered unemployed & \# of new contracts \\
\midrule
\multicolumn{3}{l}{\textbf{Panel A. Adjustment period}} \\
AI exposure $\times$ $\mathbb{1}\{F_i < \widetilde F\}$ & 0.043$^{***}$ & 0.010 \\
 & (0.009) & (0.010)\\
\quad Impact of a 10 pp increase (percent) & 4.3 & 1.0 \\
AI exposure $\times$ $\mathbb{1}\{F_i \geq \widetilde F\}$ & -0.004 & 0.015 \\
 & (0.006) & (0.012)\\
\quad Impact of a 10 pp increase (percent) & -0.4 & 1.5 \\
\addlinespace[2pt]
\multicolumn{3}{l}{\textbf{Panel B. Later post-treatment period}} \\
AI exposure $\times$ $\mathbb{1}\{F_i < \widetilde F\}$ & 0.077$^{***}$ & 0.008 \\
 & (0.018) & (0.014)\\
\quad Impact of a 10 pp increase (percent) & 7.7 & 0.8 \\
AI exposure $\times$ $\mathbb{1}\{F_i \geq \widetilde F\}$ & -0.013 & 0.012 \\
 & (0.011) & (0.014)\\
\quad Impact of a 10 pp increase (percent) & -1.3 & 1.2 \\
Observations & 31,122 & 30,786 \\
CNO4 FE & Yes & Yes \\
CNO1 $\times$ year-month FE & Yes & Yes \\
Clustered standard errors & CNO4 & CNO4 \\
\bottomrule
\end{tabular}
\begin{tablenotes}[flushleft]
\footnotesize
\item \emph{Notes:} Each column reports one regression: column 1 uses registered unemployment and column 2 uses new contracts. The rows report linear combinations of the exposure-by-period-by-group coefficients from that regression. The indicators equal one for occupations below the CNO2 female-employment-share median of 0.303 or at or above it, respectively. The adjustment and later periods cover event times 0--24 and 25--40. Exposure is measured in ten-percentage-point units. Standard errors in parentheses are clustered by CNO4. Entries are marginal effects, not ATTs; the impact rows report 100 times the corresponding log-point coefficient and carry no additional significance markers. $^{***}p<0.01$, $^{**}p<0.05$, and $^{*}p<0.10$.
\end{tablenotes}
\end{threeparttable}
\end{table}


Table~\ref{tab:v1_feminization_median_split} reports one regression for each outcome. Below-median occupations have unemployment gradients of 4.3 percent during adjustment and 7.7 percent later. The corresponding estimates for more feminized occupations are close to zero or negative and imprecise. New-contract estimates are small in both groups. This pattern is consistent with larger displacement or slower unemployment exits in less feminized occupations, but it is not robust enough to establish a mechanism: the subgroup pre-trend tests reject for part of the sample, and detrending materially changes the dynamic paths reported in Appendix~\ref{app:v1_feminization}.

\subsection{Further Results}
\label{sec:v1_further}

This subsection asks whether the main conclusion survives estimators that construct the counterfactual differently. Both alternatives are informative, but neither reproduces the preferred within-CNO1 comparison while retaining a continuous treatment.

\subsubsection{Continuous-DiD Alternative}
\label{sec:v1_contdid}

The public \texttt{contdid} implementation does not admit the CNO1-by-month effects used in the preferred design. We therefore report three transparent approximations instead of treating residualization as an equivalent conditional estimator. The CNO1-stratified estimator applies continuous DiD separately within families with at least 15 positive-exposure and 5 zero-exposure occupations and aggregates family-specific average causal responses using positive-exposure occupation shares. It retains CNO1 groups 1--5 and about 88 percent of positive-exposure occupations. A second estimator subtracts the change among zero-exposure occupations in each supported family before pooling. The third is the unconditional continuous-DiD comparison.

Appendix Table~\ref{tab:v1_contdid_alternatives} reports average causal response on the treated (ACRT) estimates over the same two intervals.\footnote{For a post-treatment window \(W\), the estimand is \(\widehat{\theta}_{W}=|W|^{-1}\sum_{t\in W}\widehat{\mathrm{ACRT}}_{t}\). If \(a_W\) assigns weight \(1/|W|\) to event times in \(W\), the estimated variance is \(a_W'\widehat{\Sigma}a_W\), where \(\widehat{\Sigma}\) is the influence-function covariance matrix. Family-specific estimates are aggregated using positive-exposure occupation shares, and disjoint-family covariance matrices enter with squared aggregation weights.} The CNO1-stratified unemployment response rises from 3.4 percent during adjustment to 7.1 percent later. The contract estimates are 2.1 and 5.7 percent. These effects are larger than the TWFE estimates and reinforce the direction of the unemployment result. They remain secondary because stratification cannot absorb family-specific shocks in every month and because the full-period pre-treatment diagnostics are less favorable. Appendix Figure~\ref{fig:v1_contdid_paths} reports the dynamic paths.

\subsubsection{Synthetic Difference-in-Differences}
\label{sec:v1_sdid}

Synthetic difference-in-differences provides a separate upper-tail design. Occupations with exposure above 0.1169 are treated, while every occupation at or below the cutoff remains eligible for the donor pool. Retaining middle-exposure occupations makes the donor pool more economically comparable than a zero-exposure-only design, but it changes the estimand from a marginal response to an upper-tail contrast.\footnote{Conceptually, \(\widehat{\tau}_{SDID}=T_1^{-1}\sum_{t>T_0}\left(\overline{Y}_{Tt}-\sum_j\widehat{\omega}_jY_{jt}\right)-\sum_{t\leq T_0}\widehat{\lambda}_t\left(\overline{Y}_{Tt}-\sum_j\widehat{\omega}_jY_{jt}\right)\), where \(\widehat{\omega}_j\) are donor weights and \(\widehat{\lambda}_t\) are pre-period time weights.}

Unit weights reproduce the treated pre-period path, and time weights emphasize pre-treatment months that best predict post-treatment donor outcomes. The adjusted variant removes CNO1-by-month means before constructing the synthetic comparison.\footnote{Inference uses 500 placebo reassignments. The CNO1-by-month projection is computed from outcome and family-month cells and is held fixed across placebo assignments. Event-study confidence intervals use the same number of reassignments.}

Appendix Table~\ref{tab:v1_sdid} reports a 2.7 percent unemployment effect during adjustment and 4.1 percent later; the corresponding contract estimates are near zero. Synthetic weighting improves pre-treatment fit and again points to a positive unemployment response. Its stronger fit comes at the cost of imposing an ad hoc binary treatment and estimating a different parameter. Appendix~\ref{app:v1_sdid} reports the paths, donor weights, and concentration diagnostics.

\subsection{Summary of Evidence}
\label{sec:v1_synthesis}

Across designs, more exposed occupations tend to experience a relative increase in registered unemployment after the release of ChatGPT. The preferred within-family TWFE estimates are economically modest and imprecise, while continuous DiD and synthetic DiD produce larger positive effects under different identifying assumptions. New contracts do not show a stable aggregate response, although younger workers exhibit more positive contract gradients. The age evidence points most clearly to workers ages 30--39 in the later period, but formal cross-group tests make this conclusion suggestive rather than definitive. The full-period pre-trend rejections and the absence of a CNO4 labor-force denominator prevent a strong causal claim. The evidence is therefore consistent with gradual AI-related labor-market adjustment, concentrated in some worker and occupation groups, but does not establish that AI penetration has raised unemployment throughout the Spanish labor market.

\clearpage
\appendix
\numberwithin{figure}{section}
\numberwithin{table}{section}

\section{Identification Support and Pre-Trend Diagnostics}
\label{app:v1_support}

\begin{table}[!htbp]
\centering
\caption{Pre-treatment coefficient tests under alternative clustering levels}
\label{tab:v1_pretrends}
\begin{threeparttable}
\scriptsize
\setlength{\tabcolsep}{4pt}
\begin{tabular}{llcccc}
\toprule
& & \multicolumn{2}{c}{CNO4 clustering} & \multicolumn{2}{c}{CNO3 clustering} \\
\cmidrule(lr){3-4}\cmidrule(lr){5-6}
Outcome & Window & Joint nullity & Equality & Joint nullity & Equality \\
\midrule
Registered unemployed & Full: -21 to -2 & 0.003 & 0.003 & 0.015 & 0.022 \\
Registered unemployed & Early: -21 to -10 & 0.027 & 0.021 & 0.134 & 0.102 \\
Registered unemployed & Recent: -10 to -2 & 0.367 & 0.280 & 0.575 & 0.479 \\
\addlinespace
New contracts & Full: -21 to -2 & 0.032 & 0.045 & 0.018 & 0.026 \\
New contracts & Early: -21 to -10 & 0.307 & 0.252 & 0.307 & 0.266 \\
New contracts & Recent: -10 to -2 & 0.049 & 0.035 & 0.070 & 0.048 \\
\bottomrule
\end{tabular}
\begin{tablenotes}[flushleft]
\footnotesize
\item \emph{Notes:} Entries are $p$-values from Wald tests of the pre-treatment coefficients in the preferred continuous-treatment specification, which includes CNO4 and CNO1-by-year-month fixed effects. The joint-nullity test evaluates whether all coefficients in the indicated window equal zero; the equality test evaluates whether they equal one another while allowing their common value to differ from zero. The full, early, and recent windows cover event times -21 to -2, -21 to -10, and -10 to -2, respectively. October 2022 (event time -1) is the omitted reference month. The first two columns allow arbitrary correlation within 502 CNO4 occupations, while the final two allow correlation within 170 CNO3 occupation groups.
\end{tablenotes}
\end{threeparttable}
\end{table}


\section{Robustness Checks}
\label{app:v1_twfe}

\begin{landscape}
\begin{table}[H]
\centering
\caption{Robustness checks: alternative outcomes and exposure measures}
\label{tab:v1_robustness}
\begin{threeparttable}
\scriptsize
\renewcommand{\arraystretch}{0.76}
\setlength{\tabcolsep}{3.5pt}
\begin{tabular}{lcccccc}
\toprule
& \multicolumn{3}{c}{\# of registered unemployed} & \multicolumn{3}{c}{\# of new contracts} \\
\cmidrule(lr){2-4}\cmidrule(lr){5-7}
& (1) & (2) & (3) & (4) & (5) & (6) \\
\midrule
\multicolumn{7}{l}{\textbf{Panel A. Baseline specification}} \\
AI exposure $\times$ adjustment period & 0.019$^{***}$ & 0.012$^{**}$ & 0.012$^{*}$ & 0.033$^{***}$ & 0.014 & 0.014 \\
 & (0.005) & (0.006) & (0.007) & (0.008) & (0.009) & (0.009) \\
Impact of a 10 pp increase (percent) & 1.9 & 1.2 & 1.2 & 3.3 & 1.4 & 1.4 \\
\addlinespace
AI exposure $\times$ later period & 0.038$^{***}$ & 0.017 & 0.017 & 0.029$^{***}$ & 0.013 & 0.013 \\
 & (0.010) & (0.012) & (0.013) & (0.010) & (0.011) & (0.012) \\
Impact of a 10 pp increase (percent) & 3.8 & 1.7 & 1.7 & 2.9 & 1.3 & 1.3 \\
$p$-value: equal phase effects & 0.001 & 0.486 & 0.529 & 0.488 & 0.807 & 0.841 \\
Observations & 31,122 & 31,122 & 31,122 & 30,787 & 30,786 & 30,786 \\
\addlinespace
\multicolumn{7}{l}{\textbf{Panel B. Alternative outcomes: log(Y+1)}} \\
AI exposure $\times$ adjustment period & 0.019$^{***}$ & 0.012$^{**}$ & 0.012$^{*}$ & 0.030$^{***}$ & 0.011 & 0.011 \\
 & (0.005) & (0.006) & (0.007) & (0.008) & (0.009) & (0.009) \\
Impact of a 10 pp increase (percent) & 1.9 & 1.2 & 1.2 & 3.0 & 1.1 & 1.1 \\
\addlinespace
AI exposure $\times$ later period & 0.037$^{***}$ & 0.017 & 0.017 & 0.025$^{**}$ & 0.008 & 0.008 \\
 & (0.010) & (0.012) & (0.013) & (0.011) & (0.012) & (0.013) \\
Impact of a 10 pp increase (percent) & 3.7 & 1.7 & 1.7 & 2.5 & 0.8 & 0.8 \\
$p$-value: equal phase effects & 0.001 & 0.485 & 0.529 & 0.401 & 0.690 & 0.726 \\
Observations & 31,124 & 31,124 & 31,124 & 31,124 & 31,124 & 31,124 \\
\addlinespace
\multicolumn{7}{l}{\textbf{Panel C. Alternative exposure: cosine-weighted}} \\
AI exposure $\times$ adjustment period & 0.025$^{***}$ & 0.018$^{***}$ & 0.018$^{**}$ & 0.031$^{***}$ & 0.012 & 0.012 \\
 & (0.005) & (0.007) & (0.008) & (0.009) & (0.010) & (0.010) \\
Impact of a 10 pp increase (percent) & 2.5 & 1.8 & 1.8 & 3.1 & 1.2 & 1.2 \\
\addlinespace
AI exposure $\times$ later period & 0.049$^{***}$ & 0.029$^{**}$ & 0.029$^{*}$ & 0.024$^{**}$ & 0.007 & 0.007 \\
 & (0.011) & (0.013) & (0.015) & (0.011) & (0.012) & (0.013) \\
Impact of a 10 pp increase (percent) & 4.9 & 2.9 & 2.9 & 2.4 & 0.7 & 0.7 \\
$p$-value: equal phase effects & 0.000 & 0.177 & 0.227 & 0.258 & 0.491 & 0.578 \\
Observations & 31,122 & 31,122 & 31,122 & 30,787 & 30,786 & 30,786 \\
\addlinespace
\midrule
CNO1 $\times$ year-month FE & No & Yes & Yes & No & Yes & Yes \\
Clustered standard errors & CNO4 & CNO4 & CNO3 & CNO4 & CNO4 & CNO3 \\
\bottomrule
\end{tabular}
\begin{tablenotes}[flushleft]
\tiny
\item \emph{Notes:} Each panel reports marginal effects for the adjustment period (event times 0--24) and the later period (event times 25--40), relative to all pre-treatment months. Panel A reproduces the baseline specification. Panel B replaces the logarithmic outcomes with $\log(Y+1)$. Panel C uses the cosine-weighted exposure measure. Columns 1 and 4 include CNO4 and year-month fixed effects; the remaining columns include CNO4 and CNO1-by-year-month fixed effects. Exposure measures are divided by 0.10. Standard errors are clustered as indicated. Impact rows report 100 times the coefficients; equality rows test whether adjustment- and later-period effects are equal. $^{***}p<0.01$, $^{**}p<0.05$, and $^{*}p<0.10$.
\end{tablenotes}
\end{threeparttable}
\end{table}
\end{landscape}

\begin{landscape}
\begin{table}[H]
\centering
\caption{Pre-treatment diagnostics for alternative outcomes and exposure measures}
\label{tab:v1_robustness_pretrends}
\begin{threeparttable}
\scriptsize
\renewcommand{\arraystretch}{0.88}
\setlength{\tabcolsep}{3.0pt}
\begin{tabular}{@{}llcccccc@{}}
\toprule
& & \multicolumn{2}{c}{Benchmark: CNO4 clustering} & \multicolumn{2}{c}{CNO1-by-month: CNO4 clustering} & \multicolumn{2}{c}{CNO1-by-month: CNO3 clustering} \\
\cmidrule(lr){3-4}\cmidrule(lr){5-6}\cmidrule(lr){7-8}
Outcome & Window & Nullity & Equality & Nullity & Equality & Nullity & Equality \\
\midrule
\multicolumn{8}{l}{\textbf{Panel A. Baseline specification}} \\
Registered unemployed & Full: -21 to -2 & $<0.001$ & $<0.001$ & 0.003 & 0.003 & 0.015 & 0.022 \\
Registered unemployed & Early: -21 to -10 & $<0.001$ & $<0.001$ & 0.027 & 0.021 & 0.134 & 0.102 \\
Registered unemployed & Recent: -10 to -2 & $<0.001$ & $<0.001$ & 0.367 & 0.280 & 0.575 & 0.479 \\
New contracts & Full: -21 to -2 & $<0.001$ & $<0.001$ & 0.032 & 0.045 & 0.018 & 0.026 \\
New contracts & Early: -21 to -10 & $<0.001$ & $<0.001$ & 0.307 & 0.252 & 0.307 & 0.266 \\
New contracts & Recent: -10 to -2 & $<0.001$ & $<0.001$ & 0.049 & 0.035 & 0.070 & 0.048 \\
\addlinespace
\multicolumn{8}{l}{\textbf{Panel B. Alternative outcomes: log(Y+1)}} \\
Registered unemployed & Full: -21 to -2 & $<0.001$ & $<0.001$ & 0.002 & 0.003 & 0.013 & 0.022 \\
Registered unemployed & Early: -21 to -10 & $<0.001$ & $<0.001$ & 0.024 & 0.019 & 0.122 & 0.095 \\
Registered unemployed & Recent: -10 to -2 & $<0.001$ & $<0.001$ & 0.369 & 0.282 & 0.570 & 0.474 \\
New contracts & Full: -21 to -2 & $<0.001$ & $<0.001$ & 0.007 & 0.009 & 0.006 & 0.007 \\
New contracts & Early: -21 to -10 & $<0.001$ & $<0.001$ & 0.028 & 0.018 & 0.025 & 0.017 \\
New contracts & Recent: -10 to -2 & $<0.001$ & $<0.001$ & 0.041 & 0.026 & 0.044 & 0.027 \\
\addlinespace
\multicolumn{8}{l}{\textbf{Panel C. Alternative exposure: cosine-weighted}} \\
Registered unemployed & Full: -21 to -2 & $<0.001$ & $<0.001$ & 0.002 & 0.004 & 0.007 & 0.027 \\
Registered unemployed & Early: -21 to -10 & $<0.001$ & $<0.001$ & 0.010 & 0.010 & 0.176 & 0.162 \\
Registered unemployed & Recent: -10 to -2 & $<0.001$ & $<0.001$ & 0.165 & 0.126 & 0.298 & 0.272 \\
New contracts & Full: -21 to -2 & $<0.001$ & $<0.001$ & 0.008 & 0.005 & 0.002 & 0.002 \\
New contracts & Early: -21 to -10 & $<0.001$ & $<0.001$ & 0.110 & 0.102 & 0.098 & 0.082 \\
New contracts & Recent: -10 to -2 & $<0.001$ & $<0.001$ & 0.019 & 0.011 & 0.043 & 0.027 \\
\addlinespace
\multicolumn{8}{l}{\textbf{Panel D. Alternative exposure: RF-relative}} \\
Registered unemployed & Full: -21 to -2 & $<0.001$ & $<0.001$ & $<0.001$ & 0.002 & $<0.001$ & 0.006 \\
Registered unemployed & Early: -21 to -10 & $<0.001$ & $<0.001$ & 0.004 & 0.003 & 0.013 & 0.010 \\
Registered unemployed & Recent: -10 to -2 & $<0.001$ & $<0.001$ & 0.088 & 0.073 & 0.199 & 0.204 \\
New contracts & Full: -21 to -2 & $<0.001$ & $<0.001$ & 0.009 & 0.008 & 0.002 & 0.004 \\
New contracts & Early: -21 to -10 & $<0.001$ & $<0.001$ & 0.016 & 0.017 & 0.003 & 0.005 \\
New contracts & Recent: -10 to -2 & $<0.001$ & $<0.001$ & 0.004 & 0.002 & 0.033 & 0.022 \\
\bottomrule
\end{tabular}
\begin{tablenotes}[flushleft]
\scriptsize
\item \emph{Notes:} Entries are $p$-values from Wald tests of the pre-treatment event-study coefficients for the specifications reported in Table~\ref{tab:v1_robustness}. The nullity test evaluates whether all coefficients in the indicated window equal zero. The equality test evaluates whether they equal one another while allowing their common value to differ from zero. The benchmark includes CNO4 and year-month fixed effects. The other specifications include CNO4 and CNO1-by-year-month fixed effects, with standard errors clustered at the indicated occupational level. The full, early, and recent windows cover event times -21 to -2, -21 to -10, and -10 to -2; October 2022 (event time -1) is the omitted reference month.
\end{tablenotes}
\end{threeparttable}
\end{table}
\end{landscape}

\begin{landscape}
\begin{table}[H]
\centering
\caption{Robustness to Jev-imputed AI exposure measures}
\label{tab:v1_jev_robustness}
\begin{threeparttable}
\scriptsize
\setlength{\tabcolsep}{3.5pt}
\begin{tabular}{lcccccc}
\toprule
& \multicolumn{3}{c}{\# of registered unemployed} & \multicolumn{3}{c}{\# of new contracts} \\
\cmidrule(lr){2-4}\cmidrule(lr){5-7}
& (1) & (2) & (3) & (4) & (5) & (6) \\
\midrule
\multicolumn{7}{l}{\textit{Panel A. Baseline specification}} \\
AI exposure & 0.035*** & 0.014 & 0.014 & 0.003 & -0.002 & -0.002 \\
 & (0.010) & (0.012) & (0.013) & (0.010) & (0.011) & (0.015) \\
Impact of a 10 pp increase (percent) & 3.5 & 1.4 & 1.4 & 0.3 & -0.2 & -0.2 \\
Observations & 502 & 502 & 502 & 492 & 492 & 492 \\
\addlinespace
\multicolumn{7}{l}{\textit{Panel B. Jev: highest-probability U.S. occupation}} \\
AI exposure & 0.044*** & 0.027** & 0.027** & -0.006 & -0.005 & -0.005 \\
 & (0.010) & (0.012) & (0.013) & (0.009) & (0.011) & (0.014) \\
Impact of a 10 pp increase (percent) & 4.4 & 2.7 & 2.7 & -0.6 & -0.5 & -0.5 \\
Observations & 502 & 502 & 502 & 492 & 492 & 492 \\
\addlinespace
\multicolumn{7}{l}{\textit{Panel C. Jev: probability-weighted U.S. occupations}} \\
AI exposure & 0.052*** & 0.033*** & 0.033** & -0.012 & -0.016 & -0.016 \\
 & (0.010) & (0.012) & (0.014) & (0.012) & (0.012) & (0.015) \\
Impact of a 10 pp increase (percent) & 5.2 & 3.3 & 3.3 & -1.2 & -1.6 & -1.6 \\
Observations & 502 & 502 & 502 & 492 & 492 & 492 \\
\addlinespace
\multicolumn{7}{l}{\textit{Panel D. Jev: directly imputed observed exposure}} \\
AI exposure & 0.084*** & 0.081*** & 0.081*** & -0.004 & -0.023 & -0.023 \\
 & (0.010) & (0.014) & (0.018) & (0.013) & (0.019) & (0.021) \\
Impact of a 10 pp increase (percent) & 8.4 & 8.1 & 8.1 & -0.4 & -2.3 & -2.3 \\
Observations & 502 & 502 & 502 & 492 & 492 & 492 \\
\addlinespace
\midrule
CNO1 fixed effects & No & Yes & Yes & No & Yes & Yes \\
Clustered standard errors & CNO4 & CNO4 & CNO3 & CNO4 & CNO4 & CNO3 \\
\bottomrule
\end{tabular}
\begin{tablenotes}[flushleft]
\footnotesize
\item \emph{Notes:} Entries are occupation-level long-difference estimates between November 2022 and November 2025. Panel A uses the Anthropic nearest-neighbor exposure. Panels B--D use, respectively, the Jev-assigned U.S. occupation with the highest probability, the probability-weighted average across U.S. occupations, and Jev's direct score for observed exposure. Columns 1 and 4 are unconditional first-difference regressions; columns 2, 3, 5, and 6 absorb CNO1 fixed effects. Standard errors are clustered as indicated. Each exposure measure is divided by 0.10, so coefficients correspond to a 10 percentage-point increase. $^{***}p<0.01$, $^{**}p<0.05$, and $^{*}p<0.10$.
\end{tablenotes}
\end{threeparttable}
\end{table}
\end{landscape}


\begin{figure}[H]
    \centering
    \caption{Correlation across AI exposure measures}
    \label{fig:v1_exposure_correlations}
    \includegraphics[width=0.84\textwidth]{\figdir/Exposure_correlation_matrix.png}
    \begin{minipage}{0.94\textwidth}
    \footnotesize \emph{Notes:} Cells report pairwise Spearman rank correlations
    across CNO4 occupations; the number of occupations used appears below each
    coefficient. Correlations use pairwise-complete observations. The BLS
    category cells remain blank until the source workbook is supplied.
    \end{minipage}
\end{figure}

\begin{figure}[H]
    \centering
    \caption{Robustness event studies: registered unemployment}
    \label{fig:v1_robust_u}
    \begin{subfigure}{0.48\textwidth}
        \includegraphics[width=\textwidth]{\figdir/Robustness_unemployed_panel1.png}
        \caption{Baseline}
    \end{subfigure}
    \hfill
    \begin{subfigure}{0.48\textwidth}
        \includegraphics[width=\textwidth]{\figdir/Robustness_unemployed_panel2.png}
        \caption{\(\log(Y+1)\)}
    \end{subfigure}

    \vspace{0.22cm}
    \begin{subfigure}{0.48\textwidth}
        \includegraphics[width=\textwidth]{\figdir/Robustness_unemployed_panel3.png}
        \caption{Cosine-weighted exposure}
    \end{subfigure}

    \begin{minipage}{0.94\textwidth}
    \footnotesize \emph{Notes:} The figure reports event-study estimates corresponding to Panels A--C of Table~\ref{tab:v1_robustness}. Panel (a) reproduces the preferred specification; panel (b) replaces the dependent variable with \(\log(Y+1)\); panel (c) replaces the nearest-neighbor exposure score with the cosine-weighted measure. Each regression includes CNO4 and CNO1-by-year-month fixed effects. October 2022 is the omitted reference month, and the vertical dashed line marks November 2022. Shaded areas are 95 percent confidence intervals based on standard errors clustered by CNO4.
    \end{minipage}
\end{figure}

\begin{figure}[H]
    \centering
    \caption{Robustness event studies: new contracts}
    \label{fig:v1_robust_c}
    \begin{subfigure}{0.48\textwidth}
        \includegraphics[width=\textwidth]{\figdir/Robustness_contracts_panel1.png}
        \caption{Baseline}
    \end{subfigure}
    \hfill
    \begin{subfigure}{0.48\textwidth}
        \includegraphics[width=\textwidth]{\figdir/Robustness_contracts_panel2.png}
        \caption{\(\log(Y+1)\)}
    \end{subfigure}

    \vspace{0.22cm}
    \begin{subfigure}{0.48\textwidth}
        \includegraphics[width=\textwidth]{\figdir/Robustness_contracts_panel3.png}
        \caption{Cosine-weighted exposure}
    \end{subfigure}

    \begin{minipage}{0.94\textwidth}
    \footnotesize \emph{Notes:} The figure reports the same specifications and confidence intervals as Figure~\ref{fig:v1_robust_u}, using the logarithm of newly registered contracts as the dependent variable.
    \end{minipage}
\end{figure}

\begin{figure}[H]
    \centering
    \caption{Jev exposure event studies: registered unemployment}
    \label{fig:v1_jev_robust_u}
    \begin{subfigure}{0.48\textwidth}
        \includegraphics[width=\textwidth]{\figdir/Robustness_unemployed_jev_nearest.png}
        \caption{Highest-probability U.S. occupation}
    \end{subfigure}
    \hfill
    \begin{subfigure}{0.48\textwidth}
        \includegraphics[width=\textwidth]{\figdir/Robustness_unemployed_jev_weighted.png}
        \caption{Probability-weighted exposure}
    \end{subfigure}
    \vspace{0.22cm}
    \begin{subfigure}{0.48\textwidth}
        \includegraphics[width=\textwidth]{\figdir/Robustness_unemployed_jev_direct.png}
        \caption{Direct score}
    \end{subfigure}
    \begin{minipage}{0.94\textwidth}
    \footnotesize \emph{Notes:} Each regression includes CNO4 and CNO1-by-year-month fixed effects. October 2022 is the omitted reference month, and the vertical dashed line marks November 2022. Shaded areas are 95 percent confidence intervals based on standard errors clustered by CNO4. Exposure is scaled to a 10 percentage-point change.
    \end{minipage}
\end{figure}

\begin{figure}[H]
    \centering
    \caption{Jev exposure event studies: new contracts}
    \label{fig:v1_jev_robust_c}
    \begin{subfigure}{0.48\textwidth}
        \includegraphics[width=\textwidth]{\figdir/Robustness_contracts_jev_nearest.png}
        \caption{Highest-probability U.S. occupation}
    \end{subfigure}
    \hfill
    \begin{subfigure}{0.48\textwidth}
        \includegraphics[width=\textwidth]{\figdir/Robustness_contracts_jev_weighted.png}
        \caption{Probability-weighted exposure}
    \end{subfigure}
    \vspace{0.22cm}
    \begin{subfigure}{0.48\textwidth}
        \includegraphics[width=\textwidth]{\figdir/Robustness_contracts_jev_direct.png}
        \caption{Direct score}
    \end{subfigure}
    \begin{minipage}{0.94\textwidth}
    \footnotesize \emph{Notes:} These event studies use the same specifications and confidence intervals as Figure~\ref{fig:v1_jev_robust_u}, with the logarithm of new contracts as the dependent variable.
    \end{minipage}
\end{figure}

\begin{table}[!htbp]
\centering
\caption{Binary-treatment robustness checks}
\label{tab:v1_binary_treatment}
\begin{threeparttable}
\scriptsize
\setlength{\tabcolsep}{3.5pt}
\begin{tabular}{lcccccc}
\toprule
& \multicolumn{3}{c}{\# of registered unemployed} & \multicolumn{3}{c}{\# of new contracts} \\
\cmidrule(lr){2-4}\cmidrule(lr){5-7}
& (1) & (2) & (3) & (4) & (5) & (6) \\
\midrule
Binary treatment $\times$ adjustment period & 0.064$^{***}$ & 0.039$^{**}$ & 0.039$^{*}$ & 0.092$^{***}$ & 0.021 & 0.021 \\
 & (0.015) & (0.017) & (0.020) & (0.025) & (0.025) & (0.024) \\
Impact of a 10 pp increase (percent) & 6.4 & 3.9 & 3.9 & 9.2 & 2.1 & 2.1 \\
\addlinespace
Binary treatment $\times$ later period & 0.124$^{***}$ & 0.058$^{*}$ & 0.058 & 0.081$^{**}$ & 0.013 & 0.013 \\
 & (0.027) & (0.031) & (0.036) & (0.033) & (0.034) & (0.034) \\
Impact of a 10 pp increase (percent) & 12.4 & 5.8 & 5.8 & 8.1 & 1.3 & 1.3 \\
$p$-value: equal phase effects & 0.000 & 0.307 & 0.333 & 0.571 & 0.691 & 0.735 \\
\midrule
CNO1 $\times$ year-month FE & No & Yes & Yes & No & Yes & Yes \\
Clustered standard errors & CNO4 & CNO4 & CNO3 & CNO4 & CNO4 & CNO3 \\
Observations & 31,122 & 31,122 & 31,122 & 30,787 & 30,786 & 30,786 \\
\bottomrule
\end{tabular}
\begin{tablenotes}[flushleft]
\footnotesize
\item \emph{Notes:} Treatment equals one for occupations with nearest-neighbor exposure above 0.1169, the 75th-percentile cutoff used in the binary design, and zero for occupations at or below the cutoff. The adjustment and later periods cover event times 0--24 and 25--40, respectively; all pre-treatment months form the omitted category. Columns 1 and 4 include CNO4 and year-month fixed effects; the remaining columns include CNO4 and CNO1-by-year-month fixed effects. Standard errors are clustered as indicated. Impact rows report 100 times the coefficients. Equality rows test whether the two phase effects are equal. $^{***}p<0.01$, $^{**}p<0.05$, and $^{*}p<0.10$.
\end{tablenotes}
\end{threeparttable}
\end{table}

\begin{table}[H]
\centering
\caption{Pre-treatment diagnostics for the binary-treatment design}
\label{tab:v1_binary_pretrends}
\begin{threeparttable}
\scriptsize
\setlength{\tabcolsep}{4pt}
\begin{tabular}{llcccc}
\toprule
& & \multicolumn{2}{c}{CNO4 clustering} & \multicolumn{2}{c}{CNO3 clustering} \\
\cmidrule(lr){3-4}\cmidrule(lr){5-6}
Outcome & Window & Jointly zero & Equal coefficients & Jointly zero & Equal coefficients \\
\midrule
Registered unemployed & Full: -21 to -2 & 0.019 & 0.018 & 0.006 & 0.007 \\
Registered unemployed & Early: -21 to -10 & 0.085 & 0.060 & 0.151 & 0.115 \\
Registered unemployed & Recent: -10 to -2 & 0.847 & 0.786 & 0.903 & 0.852 \\
\addlinespace
New contracts & Full: -21 to -2 & 0.256 & 0.211 & 0.196 & 0.169 \\
New contracts & Early: -21 to -10 & 0.690 & 0.662 & 0.754 & 0.715 \\
New contracts & Recent: -10 to -2 & 0.368 & 0.285 & 0.283 & 0.218 \\
\bottomrule
\end{tabular}
\begin{tablenotes}[flushleft]
\footnotesize
\item \emph{Notes:} Entries are $p$-values from Wald tests of the pre-treatment coefficients in the preferred binary-treatment specification, which includes CNO4 and CNO1-by-year-month fixed effects. Treated occupations have exposure above 0.1169, the 75th percentile of exposure; occupations at or below the cutoff form the comparison group. The joint-nullity test evaluates whether all coefficients in the indicated window equal zero, while the equality test evaluates whether they equal one another. October 2022 is the omitted reference month. Inference allows arbitrary correlation within either CNO4 occupations or CNO3 occupation groups, as indicated.
\end{tablenotes}
\end{threeparttable}
\end{table}


\clearpage
\begin{figure}[H]
    \centering
    \caption{Event-study estimates under the binary treatment definition}
    \label{fig:v1_binary_event}
    \begin{subfigure}{\panelwidth}
        \centering
        \includegraphics[width=\textwidth]{\figdir/Binary_event_unemployed.png}
        \caption{Number of registered unemployed}
    \end{subfigure}

    \vspace{0.30cm}

    \begin{subfigure}{\panelwidth}
        \centering
        \includegraphics[width=\textwidth]{\figdir/Binary_event_contracts.png}
        \caption{Number of new contracts}
    \end{subfigure}

    \begin{minipage}{0.92\textwidth}
    \footnotesize \emph{Notes:} The figure reports the preferred TWFE
    event-study specification under the binary treatment definition. Treatment
    equals one for occupations with exposure above 0.1169 and zero otherwise,
    so middle-exposure occupations remain in the control group. All
    specifications include CNO4 and CNO1-by-year-month fixed effects, with
    October 2022 as the omitted month. Shaded areas are 95 percent confidence
    intervals based on standard errors clustered by CNO4.
    \end{minipage}
\end{figure}

\begin{table}[!htbp]
\centering
\caption{Province-panel robustness}
\label{tab:v1_province}
\begin{threeparttable}
\scriptsize
\setlength{\tabcolsep}{3.5pt}
\begin{tabular}{lcccccc}
\toprule
& \multicolumn{3}{c}{\# of registered unemployed} & \multicolumn{3}{c}{\# of new contracts} \\
\cmidrule(lr){2-4}\cmidrule(lr){5-7}
& (1) & (2) & (3) & (4) & (5) & (6) \\
\midrule
AI exposure $\times$ adjustment period & 0.014$^{***}$ & 0.009$^{*}$ & 0.009 & 0.023$^{***}$ & 0.009 & 0.009 \\
 & (0.004) & (0.005) & (0.006) & (0.006) & (0.007) & (0.007) \\
Impact of a 10 pp increase (percent) & 1.4 & 0.9 & 0.9 & 2.3 & 0.9 & 0.9 \\
\addlinespace
AI exposure $\times$ later period & 0.025$^{***}$ & 0.010 & 0.010 & 0.023$^{***}$ & 0.009 & 0.009 \\
 & (0.009) & (0.011) & (0.013) & (0.008) & (0.009) & (0.009) \\
Impact of a 10 pp increase (percent) & 2.5 & 1.0 & 1.0 & 2.3 & 0.9 & 0.9 \\
$p$-value: equal phase effects & 0.068 & 0.885 & 0.900 & 0.879 & 0.940 & 0.945 \\
\midrule
CNO1 $\times$ year-month FE & No & Yes & Yes & No & Yes & Yes \\
Clustered standard errors & CNO4 & CNO4 & CNO3 & CNO4 & CNO4 & CNO3 \\
Observations & 1,329,773 & 1,329,773 & 1,329,773 & 1,033,581 & 1,033,581 & 1,033,581 \\
\bottomrule
\end{tabular}
\begin{tablenotes}[flushleft]
\footnotesize
\item \emph{Notes:} Entries are marginal effects from the province-by-CNO4 panel. The adjustment and later periods cover event times 0--24 and 25--40, respectively; all pre-treatment months form the omitted category. Every specification includes province-by-CNO4 and province-by-year-month fixed effects. Columns 2, 3, 5, and 6 additionally include CNO1-by-year-month fixed effects. Standard errors are clustered as indicated. Missing May 2024 province cells are reconstructed from June levels and published June-over-May changes when the rounded rate uniquely identifies an integer value; unresolved cells remain missing. Exposure is divided by 0.10. $^{***}p<0.01$, $^{**}p<0.05$, and $^{*}p<0.10$.
\end{tablenotes}
\end{threeparttable}
\end{table}


\clearpage
\begin{figure}[H]
    \centering
    \caption{Province-panel event-study estimates}
    \label{fig:v1_province_event}
    \begin{subfigure}{\panelwidth}
        \includegraphics[width=\textwidth]{\figdir/Province_event_unemployed.png}
        \caption{Number of registered unemployed}
    \end{subfigure}

    \vspace{0.25cm}
    \begin{subfigure}{\panelwidth}
        \includegraphics[width=\textwidth]{\figdir/Province_event_contracts.png}
        \caption{Number of new contracts}
    \end{subfigure}

    \begin{minipage}{0.92\textwidth}
    \footnotesize \emph{Notes:} The figure reports event-study estimates from
    the province-by-CNO4 panel. Both panels include province-by-CNO4,
    province-by-year-month, and CNO1-by-year-month fixed effects. Each marker
    is the marginal effect of a 10 percentage-point increase in AI exposure at
    the indicated event month. October 2022 is the omitted reference month;
    the vertical dashed line marks November 2022. Shaded areas are 95 percent
    confidence intervals based on standard errors clustered by CNO4. May 2024
    cells are recovered from June levels and published June-over-May changes
    when those values uniquely imply an integer count. Of 52,208 expected
    province-outcome cells, 2,323 remain unresolved and missing.
    \end{minipage}
\end{figure}

\begin{figure}[H]
    \centering
    \caption{Leave-one-CNO1-out sensitivity}
    \label{fig:v1_leave_one_out}
    \begin{subfigure}{0.48\textwidth}
        \centering
        \includegraphics[width=\textwidth]{\figdir/LeaveOneCNO1_unemployed.png}
        \caption{Number of registered unemployed}
    \end{subfigure}
    \hfill
    \begin{subfigure}{0.48\textwidth}
        \centering
        \includegraphics[width=\textwidth]{\figdir/LeaveOneCNO1_contracts.png}
        \caption{Number of new contracts}
    \end{subfigure}

    \begin{minipage}{0.94\textwidth}
    \footnotesize \emph{Notes:} Each marker reports the November
    2022--November 2025 long-difference estimate after omitting the indicated
    CNO1 family. The regressions absorb CNO1 fixed effects in the differenced
    equation and therefore compare occupations within the remaining CNO1
    families. Vertical bars are 95 percent confidence intervals based on
    standard errors clustered by CNO4.
    \end{minipage}
\end{figure}

\begin{figure}[H]
    \centering
    \caption{Within-family variation in AI exposure}
    \label{fig:v1_support}
    \begin{subfigure}{0.48\textwidth}
        \centering
        \includegraphics[width=\textwidth]{\figdir/Support_panelA.png}
        \caption{Within-CNO1 variation}
    \end{subfigure}
    \hfill
    \begin{subfigure}{0.48\textwidth}
        \centering
        \includegraphics[width=\textwidth]{\figdir/Support_panelB.png}
        \caption{Within-CNO2 variation}
    \end{subfigure}

    \begin{minipage}{0.94\textwidth}
    \footnotesize \emph{Notes:} The histograms report the distribution of the
    within-family standard deviation of nearest-neighbor AI exposure. Panel
    (a) groups CNO4 occupations by CNO1 family and panel (b) groups them by
    CNO2 family. The vertical dashed line marks the median family. Values near
    zero indicate limited identifying variation after the corresponding
    family-by-month fixed effects are absorbed.
    \end{minipage}
\end{figure}

\begin{table}[!htbp]
\centering
\caption{Within-family variation in AI exposure}
\label{tab:v1_family_support}
\begin{threeparttable}
\begin{tabular}{lcccc}
\toprule
Family level & Families & Median SD & Median range & Zero-range share \\
\midrule
CNO1 & 10 & 0.081 & 0.283 & 10.0\% \\
CNO2 & 62 & 0.038 & 0.083 & 33.9\% \\
\bottomrule
\end{tabular}
\begin{tablenotes}[flushleft]
\footnotesize
\item \emph{Notes:} Statistics are computed across the 502 CNO4 occupations using the nearest-neighbor exposure score. A zero range means that all CNO4 occupations inside the family have the same exposure. Smaller within-family dispersion implies weaker identifying variation after family-by-month effects are absorbed.
\end{tablenotes}
\end{threeparttable}
\end{table}


\clearpage
\begin{table}[H]
\centering
\caption{Robustness check: trailing 12-month contract registrations}
\label{tab:v1_trailing_contracts}
\begin{threeparttable}
\small
\begin{tabular}{lccc}
\toprule
& (1) & (2) & (3) \\
\midrule
AI exposure $\times$ adjustment period & 0.030$^{***}$ & 0.015 & 0.015 \\
 & (0.009) & (0.009) & (0.009) \\
Impact of a 10 pp increase (percent) & 3.0 & 1.5 & 1.5 \\
\addlinespace
AI exposure $\times$ later period & 0.026$^{**}$ & 0.009 & 0.009 \\
 & (0.011) & (0.012) & (0.013) \\
Impact of a 10 pp increase (percent) & 2.6 & 0.9 & 0.9 \\
$p$-value: equal phase effects & 0.579 & 0.475 & 0.513 \\
\midrule
CNO1 $\times$ year-month FE & No & Yes & Yes \\
Clustered standard errors & CNO4 & CNO4 & CNO3 \\
Observations & 26,069 & 26,069 & 26,069 \\
\bottomrule
\end{tabular}
\begin{tablenotes}[flushleft]
\footnotesize
\item \emph{Notes:} The dependent variable is the logarithm of contract registrations accumulated over the current and preceding 11 months. This trailing sum is a cumulative hiring flow, not a stock of active contracts. The first available observation is December 2021, leaving 11 pre-treatment months. The adjustment and later periods cover event times 0--24 and 25--40; all available pre-treatment months form the omitted category. Because adjacent outcomes share 11 monthly observations, the transformation mechanically smooths and delays the dynamic response. Exposure is measured in 10 percentage-point units. Standard errors are clustered as indicated. $^{***}p<0.01$, $^{**}p<0.05$, and $^{*}p<0.10$.
\end{tablenotes}
\end{threeparttable}
\end{table}


\begin{figure}[H]
    \centering
    \caption{Event-study estimates for trailing 12-month contract registrations}
    \label{fig:v1_trailing_contracts}
    \begin{subfigure}{\panelwidth}
        \centering
        \includegraphics[width=\textwidth]{\figdir/Trailing12_contracts_event.png}
        \caption{Number of new contracts accumulated over 12 months}
    \end{subfigure}

    \begin{minipage}{0.92\textwidth}
    \footnotesize \emph{Notes:} The dependent variable is the logarithm of
    newly registered contracts accumulated over the current and preceding 11
    months. The specification includes CNO4 and CNO1-by-year-month fixed
    effects. October 2022 is the omitted event-study month, and the vertical
    dashed line marks November 2022. Each marker is the marginal effect of a
    10 percentage-point increase in AI exposure; shaded areas are 95 percent
    confidence intervals based on standard errors clustered by CNO4. Because
    a complete trailing window is first available in December 2021, the graph
    has 11 pre-treatment months. Adjacent outcomes share 11 component months,
    so the transformation mechanically smooths and delays the response. The
    measure is a cumulative hiring flow, not a stock of active contracts.
    \end{minipage}
\end{figure}

\clearpage
\section{Heterogeneity Analysis}
\label{app:v1_heterogeneity}

\begingroup
\scriptsize
\setlength{\tabcolsep}{3.5pt}
\begin{longtable}{llcccc}
\caption{Pre-treatment diagnostics for the heterogeneity analysis}\label{tab:v1_heterogeneity_pretrends}\\
\toprule
& & \multicolumn{2}{c}{CNO4 clustering} & \multicolumn{2}{c}{CNO3 clustering} \\
\cmidrule(lr){3-4}\cmidrule(lr){5-6}
Outcome & Window & Jointly zero & Equal coefficients & Jointly zero & Equal coefficients \\
\midrule
\endfirsthead
\multicolumn{6}{c}{\tablename\ \thetable{} -- continued from previous page} \\
\toprule
Outcome & Window & Jointly zero & Equal coefficients & Jointly zero & Equal coefficients \\
\midrule
\endhead
\multicolumn{6}{l}{\textit{Age: under 30}} \\
Registered unemployed & Full: -21 to -2 & 0.307 & 0.274 & 0.135 & 0.106 \\
Registered unemployed & Early: -21 to -10 & 0.119 & 0.115 & 0.201 & 0.173 \\
Registered unemployed & Recent: -10 to -2 & 0.661 & 0.605 & 0.585 & 0.531 \\
New contracts & Full: -21 to -2 & 0.061 & 0.063 & 0.005 & 0.005 \\
New contracts & Early: -21 to -10 & 0.196 & 0.172 & 0.134 & 0.115 \\
New contracts & Recent: -10 to -2 & 0.053 & 0.056 & 0.030 & 0.036 \\
\addlinespace
\multicolumn{6}{l}{\textit{Age: 30--39}} \\
Registered unemployed & Full: -21 to -2 & 0.594 & 0.556 & 0.262 & 0.233 \\
Registered unemployed & Early: -21 to -10 & 0.369 & 0.336 & 0.172 & 0.147 \\
Registered unemployed & Recent: -10 to -2 & 0.475 & 0.464 & 0.393 & 0.461 \\
New contracts & Full: -21 to -2 & 0.128 & 0.100 & 0.228 & 0.209 \\
New contracts & Early: -21 to -10 & 0.487 & 0.475 & 0.456 & 0.449 \\
New contracts & Recent: -10 to -2 & 0.020 & 0.012 & 0.049 & 0.037 \\
\addlinespace
\multicolumn{6}{l}{\textit{Age: 40 or older}} \\
Registered unemployed & Full: -21 to -2 & 0.034 & 0.025 & 0.032 & 0.025 \\
Registered unemployed & Early: -21 to -10 & 0.047 & 0.038 & 0.034 & 0.031 \\
Registered unemployed & Recent: -10 to -2 & 0.746 & 0.658 & 0.877 & 0.818 \\
New contracts & Full: -21 to -2 & 0.077 & 0.077 & 0.016 & 0.015 \\
New contracts & Early: -21 to -10 & 0.322 & 0.257 & 0.283 & 0.228 \\
New contracts & Recent: -10 to -2 & 0.019 & 0.012 & 0.018 & 0.013 \\
\addlinespace
\multicolumn{6}{l}{\textit{Gender: men}} \\
Registered unemployed & Full: -21 to -2 & 0.135 & 0.105 & 0.093 & 0.076 \\
Registered unemployed & Early: -21 to -10 & 0.293 & 0.291 & 0.288 & 0.289 \\
Registered unemployed & Recent: -10 to -2 & 0.428 & 0.342 & 0.484 & 0.396 \\
New contracts & Full: -21 to -2 & 0.016 & 0.018 & 0.015 & 0.011 \\
New contracts & Early: -21 to -10 & 0.087 & 0.060 & 0.014 & 0.010 \\
New contracts & Recent: -10 to -2 & 0.037 & 0.031 & 0.041 & 0.028 \\
\addlinespace
\multicolumn{6}{l}{\textit{Gender: women}} \\
Registered unemployed & Full: -21 to -2 & 0.035 & 0.026 & 0.037 & 0.027 \\
Registered unemployed & Early: -21 to -10 & 0.023 & 0.015 & 0.042 & 0.034 \\
Registered unemployed & Recent: -10 to -2 & 0.856 & 0.795 & 0.933 & 0.890 \\
New contracts & Full: -21 to -2 & 0.115 & 0.089 & 0.033 & 0.026 \\
New contracts & Early: -21 to -10 & 0.098 & 0.086 & 0.062 & 0.090 \\
New contracts & Recent: -10 to -2 & 0.073 & 0.047 & 0.064 & 0.041 \\
\bottomrule
\end{longtable}
\begin{minipage}{0.94\textwidth}\footnotesize\emph{Notes:} Entries are $p$-values from Wald tests of the pre-treatment coefficients in subgroup-specific continuous-treatment regressions with CNO4 and CNO1-by-year-month fixed effects. The joint-nullity test evaluates whether all coefficients in the indicated window equal zero; the equality test evaluates whether they equal one another. October 2022 is the omitted reference month. Inference allows arbitrary correlation within CNO4 occupations or CNO3 occupation groups, as indicated.\end{minipage}
\endgroup

\begin{table}[!htbp]
\centering
\caption{Formal tests of age-group heterogeneity}
\label{tab:v1_age_cross_group_tests}
\begin{threeparttable}
\small
\setlength{\tabcolsep}{5pt}
\begin{tabular}{lcccc}
\toprule
& \multicolumn{2}{c}{\# of registered unemployed} & \multicolumn{2}{c}{\# of new contracts} \\
\cmidrule(lr){2-3}\cmidrule(lr){4-5}
& Adjustment & Later & Adjustment & Later \\
\midrule
All age-group effects equal & 0.450 & 0.105 & 0.025 & 0.010 \\
Ages 30--39 = under 30 & 0.897 & 0.320 & 0.102 & 0.274 \\
Ages 30--39 = age 40 or older & 0.232 & 0.035 & 0.050 & 0.007 \\
Under 30 = age 40 or older & 0.405 & 0.474 & 0.007 & 0.006 \\
\midrule
Observations & \multicolumn{2}{c}{92,368} & \multicolumn{2}{c}{90,829} \\
\bottomrule
\end{tabular}
\begin{tablenotes}[flushleft]
\footnotesize
\item \emph{Notes:} Entries are $p$-values from Wald tests computed in a pooled age-by-CNO4 panel. The regression includes age-by-CNO4 and age-by-CNO1-by-year-month fixed effects and permits the exposure gradient to differ across the three age groups in each post-treatment interval. Standard errors allow arbitrary correlation across age groups within CNO4 occupations.
\end{tablenotes}
\end{threeparttable}
\end{table}


\begin{figure}[H]
    \centering
    \caption{Heterogeneity in registered unemployment}
    \label{fig:v1_heterogeneity_u}
    \begin{subfigure}{0.48\textwidth}
        \includegraphics[width=\textwidth]{\figdir/Heterogeneity_unemployed_panel1.png}
        \caption{Under 30}
    \end{subfigure}
    \hfill
    \begin{subfigure}{0.48\textwidth}
        \includegraphics[width=\textwidth]{\figdir/Heterogeneity_unemployed_panel2.png}
        \caption{Ages 30--39}
    \end{subfigure}

    \vspace{0.20cm}
    \begin{subfigure}{0.48\textwidth}
        \includegraphics[width=\textwidth]{\figdir/Heterogeneity_unemployed_panel3.png}
        \caption{Age 40 or older}
    \end{subfigure}
    \hfill
    \begin{subfigure}{0.48\textwidth}
        \includegraphics[width=\textwidth]{\figdir/Heterogeneity_unemployed_panel4.png}
        \caption{Men}
    \end{subfigure}

    \vspace{0.20cm}
    \begin{subfigure}{0.48\textwidth}
        \includegraphics[width=\textwidth]{\figdir/Heterogeneity_unemployed_panel5.png}
        \caption{Women}
    \end{subfigure}

    \begin{minipage}{0.94\textwidth}
    \footnotesize \emph{Notes:} The figure reports subgroup-specific estimates from the preferred continuous-treatment specification, which includes CNO4 and CNO1-by-year-month fixed effects. Each marker is the marginal effect of a 10 percentage-point increase in exposure at the indicated event month. October 2022 is the omitted reference month and the vertical dashed line marks November 2022. Shaded areas are 95 percent confidence intervals based on standard errors clustered by CNO4. All panels use the same vertical scale. SEPE omits the age breakdown for May 2024; the age panels recover uniquely implied integer cells from June levels and the published June-over-May changes. Of 6,024 original age-outcome cells, 55 remain unresolved and missing.
    \end{minipage}
\end{figure}

\begin{figure}[H]
    \centering
    \caption{Heterogeneity in new contracts}
    \label{fig:v1_heterogeneity_c}
    \begin{subfigure}{0.48\textwidth}
        \includegraphics[width=\textwidth]{\figdir/Heterogeneity_contracts_panel1.png}
        \caption{Under 30}
    \end{subfigure}
    \hfill
    \begin{subfigure}{0.48\textwidth}
        \includegraphics[width=\textwidth]{\figdir/Heterogeneity_contracts_panel2.png}
        \caption{Ages 30--39}
    \end{subfigure}

    \vspace{0.20cm}
    \begin{subfigure}{0.48\textwidth}
        \includegraphics[width=\textwidth]{\figdir/Heterogeneity_contracts_panel3.png}
        \caption{Age 40 or older}
    \end{subfigure}
    \hfill
    \begin{subfigure}{0.48\textwidth}
        \includegraphics[width=\textwidth]{\figdir/Heterogeneity_contracts_panel4.png}
        \caption{Men}
    \end{subfigure}

    \vspace{0.20cm}
    \begin{subfigure}{0.48\textwidth}
        \includegraphics[width=\textwidth]{\figdir/Heterogeneity_contracts_panel5.png}
        \caption{Women}
    \end{subfigure}

    \begin{minipage}{0.94\textwidth}
    \footnotesize \emph{Notes:} The figure reports the same subgroup specifications and confidence intervals as Figure~\ref{fig:v1_heterogeneity_u}, using the logarithm of newly registered contracts as the dependent variable. All panels use the same vertical scale.
    \end{minipage}
\end{figure}

\section{Further Results}
\label{app:v1_further}

\subsection{Occupational Feminization}
\label{app:v1_feminization}

\begin{table}[!htbp]
\centering
\caption{AI exposure and occupational feminization}
\label{tab:v1_feminization_main}
\begin{threeparttable}
\scriptsize
\setlength{\tabcolsep}{3.5pt}
\begin{tabular}{lcccccc}
\toprule
& \multicolumn{3}{c}{\# of registered unemployed} & \multicolumn{3}{c}{\# of new contracts} \\
\cmidrule(lr){2-4}\cmidrule(lr){5-7}
& 25th pct. & Median & 75th pct. & 25th pct. & Median & 75th pct. \\
\midrule
AI exposure $\times$ adjustment period & 0.030$^{***}$ & 0.024$^{***}$ & 0.004 & 0.009 & 0.011 & 0.017$^{*}$ \\
 & (0.008) & (0.007) & (0.006) & (0.010) & (0.009) & (0.010) \\
Impact of a 10 pp increase (percent) & 3.0 & 2.4 & 0.4 & 0.9 & 1.1 & 1.7 \\
\addlinespace
AI exposure $\times$ later period & 0.054$^{***}$ & 0.042$^{***}$ & -0.000 & 0.003 & 0.006 & 0.018 \\
 & (0.015) & (0.013) & (0.011) & (0.013) & (0.012) & (0.012) \\
Impact of a 10 pp increase (percent) & 5.4 & 4.2 & -0.0 & 0.3 & 0.6 & 1.8 \\
\addlinespace
Observations & 31,122 & 31,122 & 31,122 & 30,786 & 30,786 & 30,786 \\
CNO4 FE & Yes & Yes & Yes & Yes & Yes & Yes \\
CNO1 $\times$ year-month FE & Yes & Yes & Yes & Yes & Yes & Yes \\
Clustered standard errors & CNO4 & CNO4 & CNO4 & CNO4 & CNO4 & CNO4 \\
\bottomrule
\end{tabular}
\begin{tablenotes}[flushleft]
\footnotesize
\item \emph{Notes:} Entries are marginal effects of a 10 percentage-point increase in AI exposure, evaluated at the indicated percentiles of the predetermined CNO2 female-employment share. The index is the female share of EPA employment pooled over 2017--2019; its 25th percentile, median, and 75th percentile are 0.234, 0.303, and 0.539. The adjustment and later periods cover event times 0--24 and 25--40. The specification includes CNO4 and CNO1-by-year-month fixed effects. Standard errors in parentheses are clustered by CNO4. Entries are marginal effects, not ATTs. The impact rows report 100 times the log-point coefficients and carry no additional significance markers. $^{***}p<0.01$, $^{**}p<0.05$, and $^{*}p<0.10$.
\end{tablenotes}
\end{threeparttable}
\end{table}

\begin{table}[!htbp]
\centering
\caption{Pre-treatment diagnostics for feminization specifications}
\label{tab:v1_feminization_pretrends}
\begin{threeparttable}
\small
\begin{tabular}{llcc}
\toprule
Outcome & Specification & Joint nullity & Joint equality \\
\midrule
\multicolumn{4}{l}{\textbf{Panel A. Continuous feminization}} \\
Registered unemployed & 25th percentile & 0.016 & 0.026 \\
Registered unemployed & 50th percentile & 0.004 & 0.007 \\
Registered unemployed & 75th percentile & 0.046 & 0.044 \\
New contracts & 25th percentile & 0.055 & 0.075 \\
New contracts & 50th percentile & 0.020 & 0.031 \\
New contracts & 75th percentile & 0.032 & 0.029 \\
\midrule
\multicolumn{4}{l}{\textbf{Panel B. Median split}} \\
Registered unemployed & Below median & 0.005 & 0.011 \\
Registered unemployed & At or above median & 0.221 & 0.226 \\
New contracts & Below median & 0.068 & 0.068 \\
New contracts & At or above median & 0.109 & 0.113 \\
\midrule
\multicolumn{4}{l}{\textbf{Panel C. Median split with detrended outcomes}} \\
Registered unemployed & Below median & 0.031 & 0.045 \\
Registered unemployed & At or above median & 0.084 & 0.300 \\
New contracts & Below median & 0.064 & 0.051 \\
New contracts & At or above median & 0.056 & 0.043 \\
\bottomrule
\end{tabular}
\begin{tablenotes}[flushleft]
\footnotesize
\item \emph{Notes:} Entries are $p$-values from tests of event-study coefficients over event times $-21$ through $-2$; October 2022 is the omitted month. The joint-nullity test sets all coefficients in the window equal to zero. The joint-equality test permits a common nonzero level but requires the coefficients to be equal to one another. Panel A evaluates continuous-feminization marginal effects at the indicated percentiles. Panel B reports the median-group results using the original outcomes. Panel C reports the same median-group specification after detrending each outcome using the pre-treatment regression of the outcome on month fixed effects and CNO4-specific linear trends. Inference clusters by CNO4.
\end{tablenotes}
\end{threeparttable}
\end{table}


\begin{figure}[H]
    \centering
    \caption{Heterogeneity by occupational feminization: registered unemployment}
    \label{fig:v1_feminization_u}
    \begin{subfigure}{\panelwidth}
        \includegraphics[width=\textwidth]{\figdir/Feminization_median_unemployed.png}
        \caption{Levels}
    \end{subfigure}

    \vspace{0.25cm}

    \begin{subfigure}{\panelwidth}
        \includegraphics[width=\textwidth]{\figdir/Feminization_median_detrended_unemployed.png}
        \caption{Pre-trended outcomes}
    \end{subfigure}

    \begin{minipage}{0.92\textwidth}
    \footnotesize \emph{Notes:} The figure reports marginal effects of a 10 percentage-point increase in AI exposure separately for occupations below and at or above the median CNO2 feminization index. The specifications include CNO4 fixed effects and CNO1-by-year-month effects; October 2022 is the omitted month and November 2022 is event time zero. Shaded areas are 95 percent confidence intervals based on standard errors clustered by CNO4. The lower panel repeats the exercise after subtracting fitted pre-treatment common month effects and CNO4-specific linear trends; it is a diagnostic for differential pre-trends rather than the preferred specification.
    \end{minipage}
\end{figure}

\begin{figure}[H]
    \centering
    \caption{Heterogeneity by occupational feminization: new contracts}
    \label{fig:v1_feminization_c}
    \begin{subfigure}{\panelwidth}
        \includegraphics[width=\textwidth]{\figdir/Feminization_median_contracts.png}
        \caption{Levels}
    \end{subfigure}

    \vspace{0.25cm}

    \begin{subfigure}{\panelwidth}
        \includegraphics[width=\textwidth]{\figdir/Feminization_median_detrended_contracts.png}
        \caption{Pre-trended outcomes}
    \end{subfigure}

    \begin{minipage}{0.92\textwidth}
    \footnotesize \emph{Notes:} The figure reports the corresponding marginal effects for the logarithm of newly registered contracts. All specifications, event-time conventions, confidence intervals, and the interpretation of the lower panel are the same as in Figure~\ref{fig:v1_feminization_u}.
    \end{minipage}
\end{figure}

\subsection{Continuous-DiD Alternative}
\label{app:v1_contdid}

\begin{table}[!htbp]
\centering
\caption{Continuous-DiD alternatives}
\label{tab:v1_contdid_alternatives}
\begin{threeparttable}
\scriptsize
\setlength{\tabcolsep}{3.5pt}
\begin{tabular}{lcccccc}
\toprule
& \multicolumn{3}{c}{\# of registered unemployed} & \multicolumn{3}{c}{\# of new contracts} \\
\cmidrule(lr){2-4}\cmidrule(lr){5-7}
& (1) & (2) & (3) & (4) & (5) & (6) \\
\midrule
AI exposure $\times$ adjustment period & 0.034$^{***}$ & 0.035$^{***}$ & 0.040$^{***}$ & 0.021 & 0.022 & 0.022 \\
 & (0.009) & (0.008) & (0.007) & (0.018) & (0.019) & (0.015) \\
Impact of a 10 pp increase (percent) & 3.4 & 3.5 & 4.0 & 2.1 & 2.2 & 2.2 \\
\addlinespace
AI exposure $\times$ later period & 0.071$^{***}$ & 0.072$^{***}$ & 0.109$^{***}$ & 0.057$^{**}$ & 0.054$^{**}$ & 0.059$^{***}$ \\
 & (0.015) & (0.014) & (0.012) & (0.025) & (0.026) & (0.021) \\
Impact of a 10 pp increase (percent) & 7.1 & 7.2 & 10.9 & 5.7 & 5.4 & 5.9 \\
$p$-value: equal phase effects & 0.008 & 0.002 & 0.000 & 0.004 & 0.010 & 0.002 \\
\midrule
CNO1-stratified & Yes & No & No & Yes & No & No \\
Control-based family-time adjustment & No & Yes & No & No & Yes & No \\
Unconditional pooled comparison & No & No & Yes & No & No & Yes \\
Observations & 20,727 & 20,727 & 31,563 & 19,530 & 19,530 & 30,303 \\
\bottomrule
\end{tabular}
\begin{tablenotes}[flushleft]
\footnotesize
\item \emph{Notes:} Entries average the monthly ACRT estimates over event times 0--24 and 25--40. Columns 1 and 4 estimate continuous DiD separately within supported CNO1 families and aggregate family-specific effects using positive-exposure occupation shares. Columns 2 and 5 subtract the zero-exposure outcome change within supported CNO1 families before estimation. Columns 3 and 6 are unconditional. Standard errors use linear combinations of the dynamic influence-function covariance matrices; disjoint-family matrices are combined using squared aggregation weights in the stratified estimator. Exposure is measured in 10 percentage-point units. These estimators do not reproduce the preferred CNO1-by-month TWFE specification. $^{***}p<0.01$, $^{**}p<0.05$, and $^{*}p<0.10$.
\end{tablenotes}
\end{threeparttable}
\end{table}


\begin{figure}[H]
    \centering
    \caption{CNO1-stratified continuous-DiD event-study estimates}
    \label{fig:v1_contdid_paths}
    \begin{subfigure}{\panelwidth}
        \includegraphics[width=\textwidth]{\figdir/ContDID_panel1.png}
        \caption{Number of registered unemployed}
    \end{subfigure}

    \vspace{0.30cm}

    \begin{subfigure}{\panelwidth}
        \includegraphics[width=\textwidth]{\figdir/ContDID_panel2.png}
        \caption{Number of new contracts}
    \end{subfigure}

    \begin{minipage}{0.92\textwidth}
    \footnotesize \emph{Notes:} The figure reports continuous-DiD estimates obtained separately within CNO1 families containing at least 15 positive-exposure and 5 zero-exposure occupations. Family-specific marginal responses are aggregated using each family's share of positive-exposure occupations. November 2022 is event time 0 and October 2022 is the reference month. Shaded areas are 95 percent confidence intervals based on the aggregated influence-function covariance matrices.
    \end{minipage}
\end{figure}

\subsection{Synthetic Difference-in-Differences}
\label{app:v1_sdid}

\begin{table}[H]
\centering
\caption{Synthetic difference-in-differences estimates}
\label{tab:v1_sdid}
\begin{threeparttable}
\small
\setlength{\tabcolsep}{4pt}
\begin{tabular}{lcccc}
\toprule
& \multicolumn{2}{c}{\# of registered unemployed} & \multicolumn{2}{c}{\# of new contracts} \\
\cmidrule(lr){2-3}\cmidrule(lr){4-5}
& (1) & (2) & (3) & (4) \\
\midrule
High exposure: adjustment period & 0.027$^{**}$ & 0.029$^{**}$ & -0.000 & 0.024 \\
 & (0.011) & (0.013) & (0.022) & (0.030) \\
Impact (percent) & 2.7 & 2.9 & -0.0 & 2.5 \\
\addlinespace
High exposure: later period & 0.041$^{*}$ & 0.019 & 0.008 & 0.029 \\
 & (0.024) & (0.030) & (0.029) & (0.037) \\
Impact (percent) & 4.1 & 1.9 & 0.8 & 3.0 \\
\midrule
Exposure threshold for treatment status & 0.1169 & 0.2 & 0.1169 & 0.2 \\
All lower-exposure occupations eligible as donors & Yes & No & Yes & No \\
Donor pool restricted to zero exposure & No & Yes & No & Yes \\
CNO1 $\times$ month residualization & Yes & Yes & Yes & Yes \\
Treated occupations: adjustment period & 124 & 84 & 118 & 82 \\
Donor occupations: adjustment period & 377 & 252 & 365 & 246 \\
Observations: adjustment period & 23,547 & 15,792 & 22,701 & 15,416 \\
Treated occupations: later period & 124 & 84 & 118 & 82 \\
Donor occupations: later period & 377 & 252 & 363 & 245 \\
Observations: later period & 19,038 & 12,768 & 18,278 & 12,426 \\
\bottomrule
\end{tabular}
\begin{tablenotes}[flushleft]
\footnotesize
\item \emph{Notes:} Columns 1 and 3 report the baseline SDID design, where treated occupations have nearest-neighbor exposure above 0.1169 and every occupation at or below that cutoff remains eligible for the donor pool. Columns 2 and 4 report the more dichotomous robustness design, where treated occupations have exposure above 0.2 and donors are restricted to zero-exposure occupations. The adjustment-period estimates use all pre-treatment months and event times 0--24. The later-period estimates use all pre-treatment months and event times 25--40; event times 0--24 are omitted from those fits. Unit and time weights are re-estimated separately for each specification, outcome, and interval. All columns residualize outcomes on CNO1-by-month indicators before constructing the synthetic comparison. Because this residualization depends only on outcomes and family-month cells, it is held fixed across placebo assignments. Standard errors in parentheses use placebo inference with 500 repetitions. Impact rows report $100[\exp(\widehat{\tau})-1]$. $^{***}p<0.01$, $^{**}p<0.05$, and $^{*}p<0.10$.
\end{tablenotes}
\end{threeparttable}
\end{table}


\begin{figure}[H]
    \centering
    \caption{CNO1-by-month-adjusted synthetic-DiD estimates}
    \label{fig:v1_sdid_projected}
    \textit{\# of registered unemployed}\par\vspace{0.10cm}
    \begin{subfigure}{0.48\textwidth}
        \includegraphics[width=\textwidth]{\figdir/SDID_adjusted_path_unemployed.png}
        \caption{Evolution}
    \end{subfigure}
    \hfill
    \begin{subfigure}{0.48\textwidth}
        \includegraphics[width=\textwidth]{\figdir/SDID_adjusted_event_unemployed.png}
        \caption{Event study}
    \end{subfigure}

    \vspace{0.25cm}
    \textit{\# of new contracts}\par\vspace{0.10cm}
    \begin{subfigure}{0.48\textwidth}
        \includegraphics[width=\textwidth]{\figdir/SDID_adjusted_path_contracts.png}
        \caption{Evolution}
    \end{subfigure}
    \hfill
    \begin{subfigure}{0.48\textwidth}
        \includegraphics[width=\textwidth]{\figdir/SDID_adjusted_event_contracts.png}
        \caption{Event study}
    \end{subfigure}

    \begin{minipage}{0.94\textwidth}
    \footnotesize \emph{Notes:} Treated occupations have nearest-neighbor AI
    exposure above 0.1169, while occupations at or below the cutoff form the
    donor pool. Before constructing the synthetic comparison, outcomes are
    adjust for CNO1-by-month indicators using the projected covariate-adjustment
    procedure implemented within synthetic DID. Panels
    (a) and (c) compare the treated outcome path with the weighted synthetic
    counterfactual; panels (b) and (d) report the corresponding event-study
    estimates. The counterfactual paths are reported as returned by the
    estimator, without an additional graphical level shift. Unit weights minimize pre-treatment discrepancies between the
    treated and donor paths, and time weights emphasize pre-treatment months
    that best predict post-treatment donor outcomes. Event-study confidence
    intervals use placebo inference with 500 repetitions.
    \end{minipage}
\end{figure}

\begin{table}[!htbp]
\centering
\caption{Synthetic-DiD donor-weight diagnostics}
\label{tab:v1_sdid_weight_diagnostics}
\begin{threeparttable}
\small
\begin{tabular}{lcc}
\toprule
& \# of registered unemployed & \# of new contracts \\
\midrule
Positive-weight donors & 377 & 363 \\
Effective number of donors & 374.6 & 362.2 \\
Largest donor weight & 0.004 & 0.003 \\
Share held by top 10 donors & 0.034 & 0.031 \\
Mean treated exposure & 0.289 & 0.293 \\
Donor-weighted exposure & 0.017 & 0.017 \\
Exposure contrast & 0.272 & 0.277 \\
\bottomrule
\end{tabular}
\begin{tablenotes}[flushleft]
\footnotesize
\item \emph{Notes:} The effective number of donors is \(1/\sum_j\widehat{\omega}_j^2\). The exposure contrast is the treated mean minus the donor-weighted mean. Calculations use the CNO1-by-month-adjusted expanded-donor synthetic-DiD specification.
\end{tablenotes}
\end{threeparttable}
\end{table}

\small
\begin{longtable}{p{0.08\textwidth}>{\raggedright\arraybackslash}p{0.43\textwidth}rrr}
\caption{Treated occupations and positive-weight synthetic-DiD donors}\label{tab:v1_sdid_donors}\\
\toprule
CNO4 & Occupation & Exposure & \shortstack{Weight: \\ unemployed} & \shortstack{Weight: \\ contracts} \\
\midrule
\endfirsthead
\multicolumn{5}{c}{\tablename\ \thetable\ -- continued} \\
\toprule
CNO4 & Occupation & Exposure & \shortstack{Weight: \\ unemployed} & \shortstack{Weight: \\ contracts} \\
\midrule
\endhead
\multicolumn{5}{l}{\textit{Panel A. Treated occupations}} \\
3820 & Computer programmers & 0.745 & & \\
4301 & Data recorders & 0.671 & & \\
3510 & Agents and commercial representatives & 0.628 & & \\
2613 & Financial analysts & 0.572 & & \\
2719 & Software and multimedia analysts and designers not classified under other headings & 0.519 & & \\
4113 & Statistical, financial and banking services office employees & 0.510 & & \\
2729 & Database and computer network specialists not classified elsewhere & 0.486 & & \\
3814 & Web technicians & 0.480 & & \\
3812 & Information technology user assistance technicians & 0.469 & & \\
3535 & Spokespersons and public relations agents & 0.453 & & \\
2652 & Public relations professionals & 0.453 & & \\
3613 & Management and administrative assistants & 0.453 & & \\
4309 & Administrative employees without public service tasks not classified under other headings & 0.453 & & \\
4500 & Administrative employees with public service tasks not classified under other headings & 0.453 & & \\
2822 & Philosophers, historians and political science professionals & 0.452 & & \\
3401 & Support professionals and exchange, stock and finance intermediaries & 0.441 & & \\
4412 & Receptionists (except hotels) & 0.434 & & \\
4411 & User information employees & 0.434 & & \\
4430 & Survey agents & 0.432 & & \\
2923 & Philologists, interpreters and translators & 0.430 & & \\
5823 & Tourist escorts & 0.405 & & \\
4421 & Travel agency employees & 0.405 & & \\
2630 & Technicians of companies and tourist activities & 0.405 & & \\
5825 & Tour guides & 0.405 & & \\
4223 & Personnel service employees & 0.405 & & \\
4112 & Personnel control and payroll employees & 0.405 & & \\
1327 & Branch directors of banks, financial services and insurance & 0.391 & & \\
1211 & CFOs & 0.391 & & \\
5895 & Astrologers, fortune tellers and the like & 0.385 & & \\
2411 & Physicists and astronomers & 0.385 & & \\
2821 & Sociologists, geographers, anthropologists, archaeologists and the like & 0.383 & & \\
1223 & Research and development directors & 0.378 & & \\
2484 & Graphic and multimedia designers & 0.367 & & \\
3614 & Secretaries of medical centers or clinics & 0.362 & & \\
2939 & Creative and performing artists not classified elsewhere & 0.356 & & \\
2931 & Plastic and visual arts artists & 0.356 & & \\
3739 & Other technicians and professionals supporting cultural and artistic activities & 0.356 & & \\
2612 & Financial and investment advisors & 0.350 & & \\
2611 & Accounting specialists & 0.348 & & \\
3811 & Computer systems operations technicians & 0.337 & & \\
2722 & System and network administrators & 0.337 & & \\
2721 & Database designers and administrators & 0.332 & & \\
5220 & Salespeople in stores and warehouses & 0.322 & & \\
5300 & Store Owner Merchants & 0.322 & & \\
3521 & Insurance mediators and agents & 0.319 & & \\
3160 & Quality control technicians in physical, chemical and engineering sciences & 0.315 & & \\
3129 & Other technicians in the physical, chemical, environmental and engineering sciences & 0.315 & & \\
3127 & Laboratory technicians and analysts in industrial chemistry & 0.315 & & \\
3121 & Technicians in physical and chemical sciences & 0.315 & & \\
2513 & Judges and magistrates & 0.311 & & \\
3403 & bookkeepers & 0.310 & & \\
4111 & Accounting Clerks & 0.310 & & \\
3322 & Senior technicians in health documentation & 0.306 & & \\
2325 & Instructors in information technologies in non-regulated education & 0.304 & & \\
2321 & Specialists in teaching and pedagogical methods & 0.304 & & \\
4445 & Bill, debt collectors and related employees & 0.299 & & \\
2230 & Secondary education teachers (except specific vocational training subjects) & 0.290 & & \\
3813 & Network technicians & 0.287 & & \\
5420 & Telemarketing operators & 0.285 & & \\
3534 & Real estate agents and managers & 0.283 & & \\
2625 & Personnel training specialists & 0.279 & & \\
2711 & Systems analysts & 0.276 & & \\
2712 & Software analysts and designers & 0.276 & & \\
2640 & Technical and medical sales professionals (except ICT) & 0.271 & & \\
2653 & Professionals in the sale of information and communications technologies & 0.271 & & \\
5210 & Store and warehouse section heads & 0.263 & & \\
1432 & Directors and managers of retail companies & 0.263 & & \\
2329 & Teachers and teaching professionals not classified under other headings & 0.262 & & \\
2413 & Chemicals & 0.261 & & \\
2713 & Analysts, programmers and web and multimedia designers & 0.249 & & \\
2921 & Writers & 0.246 & & \\
3323 & Higher dietetics technicians & 0.245 & & \\
2621 & Management and organization analysts & 0.243 & & \\
2810 & Economists & 0.242 & & \\
1222 & Advertising and public relations directors & 0.232 & & \\
3128 & Metallurgy and mining technicians & 0.229 & & \\
2324 & Non-formal arts education teachers & 0.214 & & \\
2323 & Non-regulated music and dance teaching teachers & 0.214 & & \\
2416 & Statisticians & 0.211 & & \\
3404 & Support professionals in statistical, mathematical and related services & 0.211 & & \\
2922 & Journalists & 0.210 & & \\
2427 & winemakers & 0.204 & & \\
2912 & Librarians, documentary filmmakers and the like & 0.203 & & \\
3522 & Purchasing agents & 0.200 & & \\
3312 & Imaging technicians for diagnosis & 0.199 & & \\
2723 & Computer network analysts & 0.199 & & \\
3731 & Photographers & 0.195 & & \\
5891 & Personal assistants or company people & 0.188 & & \\
5833 & Building janitors & 0.188 & & \\
3402 & Commercial loans and credits & 0.186 & & \\
3611 & Secretarial supervisors & 0.186 & & \\
1324 & Senior Social Services Directors & 0.181 & & \\
1113 & Directors of social interest organizations & 0.181 & & \\
1325 & Directors of other social services & 0.181 & & \\
4222 & Printing coders and correctors & 0.178 & & \\
9420 & Advertising delivery people, shoe shiners and other street workers & 0.176 & & \\
9410 & street vendors & 0.176 & & \\
5412 & Vendors in occasional markets and flea markets & 0.176 & & \\
5499 & Sellers not classified under other headings & 0.176 & & \\
5491 & Home sellers & 0.176 & & \\
5411 & Sellers at kiosks & 0.176 & & \\
2651 & Advertising and marketing professionals & 0.173 & & \\
4422 & Hotel receptionists & 0.172 & & \\
2592 & Attorneys & 0.167 & & \\
2511 & Lawyers & 0.167 & & \\
2599 & Legal professionals not classified under other headings & 0.167 & & \\
3313 & Technicians in pathological anatomy and cytology & 0.158 & & \\
5894 & Driving school instructors & 0.156 & & \\
1321 & Directors of information and communications technology (ICT) services & 0.156 & & \\
2426 & Environmental protection professionals & 0.144 & & \\
2326 & Environmental education professionals & 0.144 & & \\
1431 & Directors and managers of wholesale trade companies & 0.138 & & \\
1509 & Directors and managers of waste management companies and other service companies not classified under other headings & 0.138 & & \\
2153 & Dietitians and nutritionists & 0.133 & & \\
5922 & Autonomous police & 0.123 & & \\
5923 & local police & 0.123 & & \\
5921 & national police & 0.123 & & \\
5840 & Workers who own small accommodation & 0.121 & & \\
1411 & Hotel directors and managers & 0.121 & & \\
1419 & Directors and managers of other hosting services companies & 0.121 & & \\
3139 & Process control technicians not classified under other headings & 0.121 & & \\
1316 & Construction company directors & 0.119 & & \\
3202 & Construction supervisors & 0.119 & & \\
3539 & Artistic and sporting representatives and other commercial service agents not classified under other headings & 0.118 & & \\
\midrule
\multicolumn{5}{l}{\textit{Panel B. Donors with positive weight in at least one outcome}} \\
9812 & Drivers of animal-drawn vehicles for the transportation of people and similar & 0.000 & 0.0043 & 0.0028 \\
2240 & Primary school teachers & 0.103 & 0.0034 & 0.0031 \\
5991 & prison guards & 0.000 & 0.0033 & 0.0000 \\
1322 & Directors of children's social services & 0.040 & 0.0033 & 0.0024 \\
5721 & Childminders in daycare centers and educational centers & 0.012 & 0.0030 & 0.0033 \\
7192 & Technical facade installers & 0.000 & 0.0033 & 0.0028 \\
2251 & Early childhood education teachers & 0.000 & 0.0033 & 0.0031 \\
2312 & Special education technical educators & 0.099 & 0.0033 & 0.0028 \\
2311 & Special education teachers & 0.099 & 0.0032 & 0.0030 \\
2322 & Non-formal language teaching teachers & 0.064 & 0.0032 & 0.0029 \\
7706 & Workers in the production of alcoholic beverages other than wine & 0.000 & 0.0032 & 0.0025 \\
4441 & Bank tellers and similar & 0.085 & 0.0032 & 0.0026 \\
7532 & Installers and repairers in electro medicine & 0.000 & 0.0031 & 0.0025 \\
2210 & Professors of universities and other higher education (except vocational training) & 0.064 & 0.0029 & 0.0031 \\
2154 & Speech therapists & 0.000 & 0.0031 & 0.0028 \\
2220 & Vocational training teachers (specific subjects) & 0.000 & 0.0030 & 0.0031 \\
8155 & Fur and leather treatment machine operators & 0.000 & 0.0026 & 0.0031 \\
2252 & Early childhood education technicians & 0.040 & 0.0031 & 0.0028 \\
7837 & Shoemakers and the like & 0.000 & 0.0027 & 0.0031 \\
3153 & Aviation pilots and related professionals & 0.000 & 0.0026 & 0.0030 \\
3621 & Support professionals of the public tax administration & 0.028 & 0.0030 & 0.0028 \\
7311 & Shapers and macheros & 0.000 & 0.0025 & 0.0030 \\
8156 & Machine operators for the manufacturing of footwear, leather goods and leather gloves & 0.000 & 0.0027 & 0.0030 \\
7836 & Tanners and leather preparers & 0.000 & 0.0026 & 0.0030 \\
7833 & Cutters for fabrics, leather, skin and other materials & 0.015 & 0.0026 & 0.0030 \\
8111 & Miners and other operators at mining facilities & 0.000 & 0.0025 & 0.0030 \\
4443 & Game room and related employees & 0.000 & 0.0025 & 0.0030 \\
3629 & Other public administration support professionals for inspection and control tasks and similar tasks & 0.000 & 0.0030 & 0.0029 \\
8311 & Locomotive drivers & 0.000 & 0.0028 & 0.0030 \\
7831 & Tailors, dressmakers, furriers and milliners & 0.034 & 0.0027 & 0.0030 \\
3339 & Other alternative therapy professionals & 0.000 & 0.0028 & 0.0030 \\
3316 & Dental prosthesis technicians & 0.000 & 0.0028 & 0.0030 \\
3204 & Chemical and pharmaceutical industry supervisors & 0.000 & 0.0027 & 0.0030 \\
3714 & Promoters of equal opportunities between women and men & 0.079 & 0.0030 & 0.0030 \\
3125 & Electronic technicians, specialty in electro medicine & 0.000 & 0.0030 & 0.0028 \\
7892 & Punchers & 0.000 & 0.0026 & 0.0030 \\
5493 & Fashion, art and advertising models & 0.000 & 0.0024 & 0.0030 \\
4424 & Teleoperators & 0.000 & 0.0027 & 0.0029 \\
7835 & Upholsterers, mattress makers and the like & 0.000 & 0.0026 & 0.0029 \\
6202 & Qualified workers in sheep and goat livestock activities & 0.000 & 0.0027 & 0.0029 \\
1421 & Restaurant directors and managers & 0.000 & 0.0025 & 0.0029 \\
3732 & Interior designers and decorators & 0.000 & 0.0027 & 0.0029 \\
3331 & Professionals of acupuncture, naturopathy, homeopathy, traditional Chinese medicine and Ayurveda & 0.000 & 0.0028 & 0.0029 \\
7834 & Hand sewers, embroiderers and the like & 0.034 & 0.0027 & 0.0029 \\
3711 & Support professionals for legal services and similar services & 0.000 & 0.0027 & 0.0029 \\
2415 & Mathematicians and actuaries & 0.054 & 0.0029 & 0.0027 \\
3734 & Chefs & 0.000 & 0.0023 & 0.0029 \\
3523 & Consignees & 0.017 & 0.0027 & 0.0029 \\
2825 & Agents of equal opportunities between women and men & 0.028 & 0.0029 & 0.0025 \\
5824 & Ground stewards & 0.048 & 0.0026 & 0.0029 \\
7112 & Structural precast erectors (concrete only) & 0.000 & 0.0029 & 0.0027 \\
1313 & Manufacturing industry directors & 0.013 & 0.0026 & 0.0029 \\
7295 & Building facade and chimney cleaning staff & 0.000 & 0.0029 & 0.0026 \\
4122 & Production Support Services Office Employees & 0.093 & 0.0026 & 0.0029 \\
7619 & Artisans not classified under other headings & 0.054 & 0.0027 & 0.0029 \\
9512 & Agricultural laborers in orchards, greenhouses, nurseries and gardens & 0.020 & 0.0027 & 0.0029 \\
2830 & Priests of different religions & 0.112 & 0.0027 & 0.0029 \\
7612 & Luthiers and similar; musical instrument tuners & 0.000 & 0.0028 & 0.0029 \\
3155 & Aeronautical safety technicians & 0.000 & 0.0029 & 0.0029 \\
7832 & Pattern makers for textile and leather products & 0.000 & 0.0027 & 0.0029 \\
3110 & Draftsmen and technical draftsmen & 0.000 & 0.0027 & 0.0029 \\
5821 & Flight attendants and waiters on airplanes, ships and trains & 0.117 & 0.0026 & 0.0029 \\
5812 & Specialists in aesthetic, wellness and related treatments & 0.000 & 0.0028 & 0.0029 \\
7899 & Officials, workers and artisans of other trades not classified under other headings & 0.054 & 0.0029 & 0.0027 \\
3531 & customs representatives & 0.017 & 0.0029 & 0.0025 \\
1212 & Human resources directors & 0.000 & 0.0027 & 0.0029 \\
9100 & Domestic employees & 0.000 & 0.0027 & 0.0029 \\
7222 & Gas fitters-installers in buildings & 0.000 & 0.0026 & 0.0029 \\
3721 & Athletes and sportsmen & 0.000 & 0.0029 & 0.0024 \\
8153 & Sewing and embroidery machine operators & 0.000 & 0.0028 & 0.0029 \\
2464 & Aeronautical technical engineers & 0.000 & 0.0026 & 0.0029 \\
3833 & Telecommunications engineering technicians & 0.000 & 0.0028 & 0.0029 \\
2412 & Meteorologists & 0.038 & 0.0021 & 0.0029 \\
3321 & Senior technicians in oral hygiene & 0.000 & 0.0029 & 0.0029 \\
4121 & Supply and inventory control employees & 0.000 & 0.0026 & 0.0029 \\
7122 & Stonemasons, cutters, carvers and stone engravers & 0.000 & 0.0026 & 0.0029 \\
3713 & Professionals supporting work and social education & 0.000 & 0.0029 & 0.0028 \\
3532 & Conference and event organizers & 0.102 & 0.0027 & 0.0029 \\
8141 & Operators of machines to manufacture rubber products and derivatives of natural resins & 0.000 & 0.0027 & 0.0029 \\
2452 & Landscape architects & 0.000 & 0.0028 & 0.0029 \\
7613 & Jewelers, goldsmiths and silversmiths & 0.000 & 0.0026 & 0.0029 \\
8154 & Textile bleaching, dyeing, printing and finishing machine operators & 0.021 & 0.0027 & 0.0029 \\
8152 & Operators of looms and other weaving machines & 0.032 & 0.0027 & 0.0029 \\
1221 & Commercial and sales directors & 0.043 & 0.0027 & 0.0029 \\
9542 & Aquaculture laborers & 0.000 & 0.0022 & 0.0029 \\
7611 & Watchmakers and precision instrument mechanics & 0.000 & 0.0029 & 0.0028 \\
1329 & Directors of other professional services companies not classified under other headings & 0.033 & 0.0028 & 0.0029 \\
9602 & Building construction laborers & 0.028 & 0.0026 & 0.0029 \\
5832 & Domestic service butlers & 0.000 & 0.0025 & 0.0029 \\
6204 & Skilled workers in beekeeping and sericulture & 0.000 & 0.0027 & 0.0029 \\
9320 & Fast food preparers & 0.000 & 0.0026 & 0.0029 \\
3733 & Technicians in art galleries, museums and libraries & 0.000 & 0.0027 & 0.0029 \\
4210 & Library and Archives Employees & 0.098 & 0.0027 & 0.0029 \\
7617 & Craftsmen in wood and similar materials & 0.054 & 0.0028 & 0.0029 \\
9434 & Meter readers and collectors of recreational and vending machines & 0.000 & 0.0024 & 0.0029 \\
6110 & Skilled workers in agricultural activities (except in orchards, greenhouses, nurseries and gardens) & 0.020 & 0.0029 & 0.0027 \\
5811 & Hairdressers & 0.030 & 0.0028 & 0.0029 \\
0012 & Non-commissioned officers of the armed forces & 0.000 & 0.0028 & 0.0000 \\
1323 & Directors-managers of health centers & 0.066 & 0.0025 & 0.0028 \\
3126 & Mechanical technicians & 0.000 & 0.0026 & 0.0028 \\
5892 & Undertakers and embalmers & 0.000 & 0.0028 & 0.0027 \\
8420 & Bus and tram drivers & 0.000 & 0.0027 & 0.0028 \\
3142 & Agricultural technicians & 0.006 & 0.0027 & 0.0028 \\
2111 & family doctors & 0.000 & 0.0028 & 0.0026 \\
2434 & Aeronautical engineers & 0.075 & 0.0022 & 0.0028 \\
7402 & Aircraft engine mechanics and adjusters & 0.000 & 0.0027 & 0.0028 \\
3612 & Legal-legal assistants & 0.000 & 0.0028 & 0.0025 \\
2624 & Specialists in personnel and related policies and services & 0.065 & 0.0027 & 0.0028 \\
3405 & Appraisers & 0.000 & 0.0028 & 0.0027 \\
3206 & Supervisors of wood and pulp and paper industries & 0.000 & 0.0026 & 0.0028 \\
3141 & Technicians in biological sciences (except in healthcare areas) & 0.060 & 0.0028 & 0.0028 \\
9310 & kitchen helpers & 0.012 & 0.0025 & 0.0028 \\
5932 & forest firefighters & 0.000 & 0.0021 & 0.0028 \\
3329 & Health technicians not classified under other headings & 0.000 & 0.0028 & 0.0026 \\
6422 & Coastal and freshwater fishermen & 0.000 & 0.0025 & 0.0028 \\
2151 & Dentists and stomatologists & 0.000 & 0.0027 & 0.0028 \\
3123 & Electrical technicians & 0.000 & 0.0026 & 0.0028 \\
2451 & Architects (except landscape architects and urban planners) & 0.078 & 0.0026 & 0.0028 \\
3724 & Monitors of recreational and entertainment activities & 0.000 & 0.0028 & 0.0028 \\
6209 & Qualified workers in livestock activities not classified under other headings & 0.000 & 0.0027 & 0.0028 \\
8199 & Operators of fixed installations and machinery not classified under other headings & 0.022 & 0.0028 & 0.0027 \\
1429 & Directors and managers of catering companies and other restaurant companies & 0.000 & 0.0028 & 0.0026 \\
7621 & Prepress process workers & 0.102 & 0.0027 & 0.0028 \\
6120 & Qualified workers in orchards, greenhouses, nurseries and gardens & 0.020 & 0.0027 & 0.0028 \\
1219 & Directors of policy and planning and other administrative departments not classified elsewhere & 0.033 & 0.0027 & 0.0028 \\
3722 & Coaches and referees of sports activities & 0.000 & 0.0028 & 0.0026 \\
5492 & Sales promoters & 0.079 & 0.0028 & 0.0028 \\
3122 & Construction technicians & 0.078 & 0.0026 & 0.0028 \\
4423 & Telephonists & 0.000 & 0.0026 & 0.0028 \\
2157 & Podiatrists & 0.000 & 0.0025 & 0.0028 \\
2466 & Mining, metallurgical and related technical engineers & 0.000 & 0.0027 & 0.0028 \\
6410 & Qualified workers in forestry and natural environment activities & 0.000 & 0.0025 & 0.0028 \\
9511 & Agricultural laborers (except in orchards, greenhouses, nurseries and gardens) & 0.020 & 0.0028 & 0.0028 \\
2437 & Environmental engineers & 0.036 & 0.0027 & 0.0028 \\
3315 & Orthoprosthetics technicians & 0.000 & 0.0023 & 0.0028 \\
1120 & CEOs and CEOs & 0.033 & 0.0026 & 0.0028 \\
2436 & Mining, metallurgical and related engineers & 0.000 & 0.0024 & 0.0028 \\
4442 & Betting sales employees & 0.000 & 0.0027 & 0.0028 \\
6203 & Qualified workers in pig farming activities & 0.000 & 0.0028 & 0.0027 \\
8132 & Machine operators for manufacturing pharmaceutical, cosmetic and related products & 0.000 & 0.0028 & 0.0027 \\
9490 & Other elementary occupations & 0.000 & 0.0028 & 0.0028 \\
7820 & Cabinetmakers and related workers & 0.000 & 0.0026 & 0.0028 \\
2155 & Opticians-optometrists & 0.000 & 0.0028 & 0.0027 \\
7893 & Product sorters and testers (except food, beverages and tobacco) & 0.032 & 0.0027 & 0.0028 \\
2935 & Actors & 0.101 & 0.0026 & 0.0028 \\
5110 & Salaried cooks & 0.012 & 0.0026 & 0.0028 \\
9433 & Delivery drivers, messengers and couriers on foot & 0.041 & 0.0026 & 0.0028 \\
2439 & Engineers not classified under other headings & 0.066 & 0.0027 & 0.0028 \\
2622 & Specialists in business policy administration & 0.033 & 0.0026 & 0.0028 \\
4123 & Logistics and passenger and merchandise transportation employees & 0.096 & 0.0026 & 0.0028 \\
2414 & Geologists and geophysicists & 0.043 & 0.0026 & 0.0028 \\
3131 & Technicians in energy production facilities & 0.026 & 0.0026 & 0.0028 \\
5722 & Child care workers at home & 0.012 & 0.0028 & 0.0028 \\
4446 & Post Office Counter Employees & 0.000 & 0.0028 & 0.0000 \\
2435 & Chemical engineers & 0.000 & 0.0027 & 0.0028 \\
2469 & Technical engineers not classified under other headings & 0.000 & 0.0028 & 0.0028 \\
3715 & Community animators & 0.000 & 0.0027 & 0.0028 \\
7131 & Carpenters (except cabinetmakers) & 0.000 & 0.0026 & 0.0028 \\
9601 & Public works laborers & 0.028 & 0.0026 & 0.0028 \\
8331 & Operators of earthmoving machinery and similar equipment & 0.000 & 0.0024 & 0.0028 \\
2512 & Prosecutors & 0.102 & 0.0026 & 0.0028 \\
3209 & Supervisors in other manufacturing industries & 0.000 & 0.0027 & 0.0028 \\
8191 & Kiln operators and glass and ceramic installations & 0.000 & 0.0028 & 0.0028 \\
8151 & Fiber preparation, spinning and winding machine operators & 0.000 & 0.0027 & 0.0028 \\
9520 & Livestock workers & 0.000 & 0.0026 & 0.0028 \\
3132 & Technicians in waste and water treatment facilities and other operators in similar plants & 0.000 & 0.0026 & 0.0028 \\
3533 & Agents or intermediaries in the hiring of labor (except entertainment representatives) & 0.000 & 0.0027 & 0.0028 \\
9222 & Vehicle cleaners & 0.000 & 0.0025 & 0.0028 \\
5629 & People care workers in health services not classified under other headings & 0.000 & 0.0028 & 0.0027 \\
8121 & Operators in facilities for obtaining and transforming metals & 0.000 & 0.0025 & 0.0028 \\
7132 & Installers of metal enclosures and metal carpenters (except assemblers of metal structures) & 0.000 & 0.0026 & 0.0028 \\
2431 & Industrial and production engineers & 0.037 & 0.0026 & 0.0028 \\
2482 & Product and clothing designers & 0.044 & 0.0027 & 0.0028 \\
2453 & Urban planners and traffic engineers & 0.096 & 0.0028 & 0.0000 \\
7522 & Power line installers and repairers & 0.000 & 0.0026 & 0.0028 \\
5710 & Home-based personal care workers & 0.000 & 0.0028 & 0.0027 \\
8142 & Machine operators for manufacturing plastic products & 0.000 & 0.0027 & 0.0028 \\
2441 & electrical engineers & 0.059 & 0.0026 & 0.0028 \\
3723 & Sports activity instructors & 0.000 & 0.0026 & 0.0028 \\
2936 & Radio, television and other presenters & 0.063 & 0.0028 & 0.0027 \\
7811 & Wood treatment workers & 0.000 & 0.0025 & 0.0028 \\
8322 & Mobile forestry machinery operators & 0.000 & 0.0022 & 0.0028 \\
9603 & Laborers in mining, quarries and other extractive industries & 0.000 & 0.0026 & 0.0028 \\
2461 & Industrial and production technical engineers & 0.000 & 0.0026 & 0.0028 \\
5831 & Maintenance and cleaning supervisors in offices, hotels and other establishments & 0.044 & 0.0027 & 0.0028 \\
2424 & Agricultural technical engineers & 0.000 & 0.0027 & 0.0028 \\
8412 & Salaried drivers of cars, taxis and vans & 0.000 & 0.0027 & 0.0028 \\
5500 & Cashiers and lockers (except banks) & 0.085 & 0.0028 & 0.0027 \\
7293 & Glaziers & 0.000 & 0.0026 & 0.0028 \\
2421 & Biologists, botanists, zoologists and the like & 0.061 & 0.0027 & 0.0028 \\
8112 & Operators in facilities for the preparation of minerals and rocks & 0.000 & 0.0026 & 0.0028 \\
3152 & Captains and deck officers & 0.000 & 0.0024 & 0.0028 \\
7705 & Fruit and vegetable canning workers and non-alcoholic beverage production workers & 0.000 & 0.0028 & 0.0026 \\
7223 & Duct installers in public works & 0.012 & 0.0026 & 0.0028 \\
6201 & Qualified workers in beef livestock activities & 0.000 & 0.0026 & 0.0028 \\
7701 & Butchers and meat industry workers & 0.000 & 0.0028 & 0.0026 \\
2465 & Chemical technical engineers & 0.000 & 0.0025 & 0.0028 \\
8440 & Motorcycle and moped drivers & 0.000 & 0.0026 & 0.0028 \\
8201 & Mechanical machinery assemblers & 0.000 & 0.0028 & 0.0027 \\
7221 & Plumbers & 0.012 & 0.0026 & 0.0028 \\
7615 & Glass blowers, shapers, laminators, cutters and polishers & 0.000 & 0.0026 & 0.0028 \\
8431 & Truck Owner Drivers & 0.000 & 0.0027 & 0.0028 \\
3203 & Food and tobacco industry supervisors & 0.056 & 0.0027 & 0.0028 \\
7709 & Food and beverage tasters and classifiers & 0.000 & 0.0028 & 0.0027 \\
8432 & Salaried truck drivers & 0.000 & 0.0027 & 0.0028 \\
7193 & Installers of waterproofing systems in buildings & 0.000 & 0.0026 & 0.0028 \\
3325 & Physiotherapist assistants & 0.000 & 0.0026 & 0.0028 \\
2481 & Technical architects and urban planning technicians & 0.078 & 0.0026 & 0.0028 \\
1112 & Management personnel of the public administration & 0.000 & 0.0025 & 0.0028 \\
8133 & Photographic and related laboratory operators & 0.011 & 0.0028 & 0.0026 \\
2443 & Telecommunications engineers & 0.000 & 0.0026 & 0.0028 \\
5931 & Firefighters (except forestry) & 0.000 & 0.0021 & 0.0028 \\
7708 & Preparers and processors of tobacco and its products & 0.000 & 0.0024 & 0.0028 \\
3205 & Supervisors of plastic, rubber and natural resin transformation industries & 0.000 & 0.0028 & 0.0026 \\
2824 & Work and social education professionals & 0.000 & 0.0028 & 0.0027 \\
7618 & Artisans in fabrics, leather and similar materials, fiber preparers and weavers with artisan looms or knitting and related fabrics & 0.054 & 0.0025 & 0.0028 \\
7623 & Binding process workers & 0.000 & 0.0027 & 0.0028 \\
1501 & Directors and managers of recreational, cultural and sports activities companies & 0.044 & 0.0028 & 0.0026 \\
2463 & Mechanical technical engineers & 0.081 & 0.0026 & 0.0028 \\
8312 & Railway shunting agents & 0.000 & 0.0028 & 0.0028 \\
2937 & Bullfighting professionals & 0.044 & 0.0021 & 0.0028 \\
6423 & Deep sea fishermen & 0.000 & 0.0023 & 0.0028 \\
8122 & Operators of metal polishing, galvanizing and coating machines & 0.000 & 0.0027 & 0.0028 \\
3326 & Technicians in occupational risk prevention and environmental health & 0.000 & 0.0027 & 0.0028 \\
2433 & mechanical engineers & 0.081 & 0.0025 & 0.0028 \\
7531 & Mechanics and repairers of electronic equipment & 0.000 & 0.0026 & 0.0028 \\
8209 & Assemblers and assemblers not classified elsewhere & 0.000 & 0.0027 & 0.0028 \\
2623 & Public administration specialists & 0.000 & 0.0028 & 0.0026 \\
5120 & Salaried waiters & 0.000 & 0.0025 & 0.0028 \\
7616 & Sign makers, glass engravers, decorative painters of various items & 0.000 & 0.0027 & 0.0028 \\
0011 & Armed forces officers & 0.000 & 0.0025 & 0.0028 \\
2473 & Telecommunications technical engineers & 0.000 & 0.0028 & 0.0027 \\
5893 & Animal caretakers and trainers & 0.000 & 0.0028 & 0.0027 \\
7191 & Building maintainers & 0.000 & 0.0026 & 0.0028 \\
2471 & Electrical technical engineers & 0.059 & 0.0028 & 0.0025 \\
7401 & Motor Vehicle Mechanics and Adjusters & 0.000 & 0.0026 & 0.0027 \\
3133 & Technicians in control of chemical product processing facilities & 0.000 & 0.0027 & 0.0027 \\
2823 & Psychologists & 0.059 & 0.0027 & 0.0027 \\
2432 & Construction and civil works engineers & 0.008 & 0.0026 & 0.0027 \\
2422 & Agronomists & 0.000 & 0.0026 & 0.0027 \\
8170 & Laundry and dry cleaning machine operators & 0.000 & 0.0027 & 0.0027 \\
5000 & Waiters and chefs-owners & 0.000 & 0.0027 & 0.0027 \\
1326 & Directors of education services & 0.052 & 0.0027 & 0.0027 \\
5942 & Security guard assistants and similar not authorized to be armed & 0.000 & 0.0027 & 0.0027 \\
2158 & Health and occupational and environmental hygiene professionals & 0.000 & 0.0026 & 0.0027 \\
8144 & Operators of sawmills, panel manufacturing machines and related installations for the treatment of wood and cork & 0.000 & 0.0026 & 0.0027 \\
7314 & Assemblers of metal structures & 0.000 & 0.0025 & 0.0027 \\
7703 & Bakers, pastry chefs and confectioners & 0.000 & 0.0027 & 0.0027 \\
3317 & Hearing aid technicians & 0.000 & 0.0026 & 0.0027 \\
9530 & Agricultural laborers & 0.000 & 0.0026 & 0.0027 \\
8340 & Deck sailors, engine sailors and the like & 0.000 & 0.0024 & 0.0027 \\
7121 & Masons & 0.000 & 0.0026 & 0.0027 \\
7199 & Other workers in structural construction works not classified under other headings & 0.000 & 0.0026 & 0.0027 \\
3831 & Audiovisual recording technicians & 0.000 & 0.0027 & 0.0027 \\
9431 & Ordinances & 0.075 & 0.0027 & 0.0027 \\
1422 & Directors and managers of bars, cafes and similar & 0.000 & 0.0027 & 0.0027 \\
3124 & Electronic technicians (except electro medicine) & 0.020 & 0.0027 & 0.0027 \\
7533 & Installers and repairers in information and communications technologies & 0.000 & 0.0027 & 0.0027 \\
8193 & Packaging, bottling and labeling machine operators & 0.000 & 0.0027 & 0.0027 \\
9229 & Other cleaning staff & 0.000 & 0.0024 & 0.0027 \\
2442 & electronic engineers & 0.100 & 0.0024 & 0.0027 \\
5621 & Assistant pharmacy technicians & 0.073 & 0.0027 & 0.0027 \\
7614 & Ceramic workers, potters and the like & 0.000 & 0.0027 & 0.0027 \\
3134 & Oil and natural gas refinery technicians & 0.041 & 0.0024 & 0.0027 \\
2911 & Archivists and museum curators & 0.000 & 0.0027 & 0.0027 \\
2423 & forestry engineers & 0.000 & 0.0024 & 0.0027 \\
3327 & Veterinary assistants & 0.018 & 0.0027 & 0.0027 \\
7111 & Formworkers and concrete construction workers & 0.000 & 0.0026 & 0.0027 \\
5612 & Primary care nursing assistants & 0.000 & 0.0027 & 0.0026 \\
6300 & Qualified workers in mixed agricultural activities & 0.000 & 0.0027 & 0.0027 \\
9443 & Sweepers and the like & 0.000 & 0.0027 & 0.0027 \\
7510 & Construction electricians and related & 0.000 & 0.0026 & 0.0027 \\
6205 & Skilled workers in poultry and rabbit farming & 0.000 & 0.0027 & 0.0027 \\
7521 & Mechanics and electrical equipment repairers & 0.000 & 0.0026 & 0.0027 \\
7250 & Refrigeration and air conditioning mechanics-installers & 0.019 & 0.0025 & 0.0027 \\
7321 & Blacksmiths and forgers & 0.000 & 0.0025 & 0.0027 \\
8159 & Operators of machines for manufacturing textile products not classified under other headings & 0.032 & 0.0027 & 0.0027 \\
9223 & window cleaners & 0.000 & 0.0026 & 0.0027 \\
5611 & Hospital nursing assistants & 0.000 & 0.0027 & 0.0027 \\
7292 & Installers of thermal insulation and soundproofing material & 0.044 & 0.0027 & 0.0027 \\
7404 & Mechanics and adjusters of naval and railway machinery & 0.000 & 0.0025 & 0.0027 \\
3324 & Optometry technicians & 0.000 & 0.0027 & 0.0027 \\
2454 & Geographic engineers and cartographers & 0.080 & 0.0024 & 0.0027 \\
2472 & Electronic technical engineers & 0.020 & 0.0026 & 0.0027 \\
7403 & Mechanics and adjusters of agricultural and industrial machinery & 0.000 & 0.0025 & 0.0027 \\
9543 & Forestry and hunting laborers & 0.000 & 0.0025 & 0.0027 \\
7291 & Cover fitters & 0.016 & 0.0025 & 0.0027 \\
3712 & Private detectives & 0.000 & 0.0027 & 0.0000 \\
5941 & Security guards and similar authorized to be armed & 0.000 & 0.0027 & 0.0027 \\
3622 & Support professionals of the public administration of social services & 0.000 & 0.0027 & 0.0026 \\
9820 & Replenishers & 0.000 & 0.0026 & 0.0027 \\
5899 & Personal service workers not classified under other headings & 0.041 & 0.0027 & 0.0027 \\
7240 & Tilers, parquet installers and the like & 0.000 & 0.0026 & 0.0027 \\
9210 & Cleaning staff for offices, hotels and other similar establishments & 0.000 & 0.0027 & 0.0027 \\
8192 & Boiler and steam engine operators & 0.000 & 0.0024 & 0.0027 \\
7704 & Workers in the processing of milk and production of dairy products (including ice cream) & 0.000 & 0.0024 & 0.0027 \\
2159 & Health professionals not classified under other headings & 0.000 & 0.0027 & 0.0027 \\
7232 & Painters in manufacturing industries & 0.000 & 0.0026 & 0.0027 \\
7231 & Painters and wallpaperers & 0.000 & 0.0026 & 0.0027 \\
3311 & Radiotherapy technicians & 0.000 & 0.0024 & 0.0027 \\
8114 & Machinery operators to manufacture products derived from non-metallic minerals & 0.000 & 0.0025 & 0.0027 \\
8131 & Operators in chemical industrial plants & 0.000 & 0.0026 & 0.0027 \\
7812 & Woodworking machine fitters and operators & 0.000 & 0.0026 & 0.0027 \\
2933 & Choreographers and dancers & 0.080 & 0.0026 & 0.0027 \\
8143 & Machine operators for manufacturing paper and cardboard products & 0.000 & 0.0027 & 0.0027 \\
3151 & Chiefs and engineering officers & 0.000 & 0.0024 & 0.0027 \\
7894 & Fumigators and other pest and weed controllers & 0.000 & 0.0027 & 0.0027 \\
8202 & Assemblers of electrical and electronic equipment & 0.000 & 0.0027 & 0.0027 \\
7212 & Paste and mortar coating applicators & 0.000 & 0.0026 & 0.0027 \\
3832 & broadcast technicians & 0.020 & 0.0026 & 0.0027 \\
9700 & Laborers in manufacturing industries & 0.000 & 0.0027 & 0.0027 \\
9441 & Waste collectors & 0.000 & 0.0025 & 0.0027 \\
7312 & Welders and oxycutters & 0.000 & 0.0025 & 0.0027 \\
7324 & Metal polishers and tool sharpeners & 0.000 & 0.0026 & 0.0027 \\
7405 & Bicycle repairers and the like & 0.000 & 0.0026 & 0.0027 \\
7294 & Solar energy panel assemblers-installers & 0.000 & 0.0026 & 0.0027 \\
5430 & Gas station dispensers & 0.000 & 0.0027 & 0.0027 \\
2934 & Film, theater and related directors & 0.092 & 0.0026 & 0.0027 \\
7313 & Sheet metal workers and boilermakers & 0.000 & 0.0025 & 0.0027 \\
1315 & Directors of supply, transportation, distribution and related companies & 0.096 & 0.0027 & 0.0027 \\
9221 & Hand dry cleaners and similar & 0.000 & 0.0027 & 0.0027 \\
2932 & Composers, musicians and singers & 0.037 & 0.0027 & 0.0026 \\
7622 & Printing process workers & 0.000 & 0.0027 & 0.0027 \\
8411 & Drivers who own cars, taxis and vans & 0.000 & 0.0027 & 0.0024 \\
9811 & Freight transport workers and unloaders & 0.000 & 0.0026 & 0.0027 \\
5993 & Forestry and environmental agents & 0.000 & 0.0025 & 0.0027 \\
8160 & Machine operators to produce food, beverage and tobacco products & 0.000 & 0.0027 & 0.0026 \\
9541 & fishing laborers & 0.000 & 0.0024 & 0.0027 \\
5999 & Protection and security services workers not classified under other headings & 0.057 & 0.0026 & 0.0027 \\
6421 & Skilled workers in aquaculture & 0.000 & 0.0026 & 0.0027 \\
5822 & Ground transportation conductors and collectors & 0.117 & 0.0027 & 0.0025 \\
2483 & Technical engineers in topography & 0.000 & 0.0024 & 0.0027 \\
8321 & Mobile agricultural machinery operators & 0.000 & 0.0025 & 0.0026 \\
5622 & Health emergency technicians & 0.000 & 0.0026 & 0.0026 \\
8145 & Operators in pulp preparation and paper manufacturing facilities & 0.000 & 0.0026 & 0.0026 \\
2130 & Veterinarians & 0.093 & 0.0026 & 0.0026 \\
7211 & plasterers & 0.000 & 0.0026 & 0.0026 \\
3201 & Supervisors in mining engineering & 0.030 & 0.0021 & 0.0026 \\
7315 & Assemblers of cabled structures and cable splicers & 0.000 & 0.0026 & 0.0026 \\
3314 & Technicians in clinical diagnostic laboratory & 0.000 & 0.0026 & 0.0026 \\
2122 & Specialized nurses (except midwives) & 0.094 & 0.0026 & 0.0025 \\
2140 & Pharmacists & 0.090 & 0.0026 & 0.0026 \\
7702 & Fish industry workers & 0.000 & 0.0025 & 0.0026 \\
3143 & Forestry and natural environment technicians & 0.022 & 0.0023 & 0.0026 \\
8333 & Forklift operators & 0.000 & 0.0026 & 0.0026 \\
2462 & Technical public works engineers & 0.008 & 0.0026 & 0.0026 \\
3207 & Production supervisors in graphic arts industries and paper products manufacturing & 0.000 & 0.0026 & 0.0026 \\
2425 & Forestry and natural environment technical engineers & 0.022 & 0.0023 & 0.0026 \\
2156 & Occupational therapists & 0.008 & 0.0025 & 0.0026 \\
7707 & Wine making workers & 0.000 & 0.0026 & 0.0026 \\
7322 & Tool manufacturing workers, mechanic-fitters, model makers, tool makers and the like & 0.000 & 0.0025 & 0.0026 \\
8113 & Surveyors and related workers & 0.000 & 0.0026 & 0.0026 \\
8332 & Operators of cranes, forklifts and similar material movement machinery & 0.000 & 0.0026 & 0.0025 \\
9432 & Baggage handlers and similar & 0.073 & 0.0026 & 0.0026 \\
9442 & Waste sorters, clean point operators and scrap collectors & 0.000 & 0.0026 & 0.0024 \\
1311 & Production directors of agricultural and forestry farms & 0.000 & 0.0024 & 0.0026 \\
5910 & civil guards & 0.000 & 0.0026 & 0.0000 \\
0020 & Troops and sailors of the armed forces & 0.000 & 0.0026 & 0.0000 \\
2152 & Physiotherapists & 0.017 & 0.0024 & 0.0026 \\
3154 & air traffic controllers & 0.000 & 0.0023 & 0.0026 \\
4221 & Postal service employees (except counter employees) & 0.000 & 0.0025 & 0.0024 \\
7323 & Machine tool adjusters and operators & 0.000 & 0.0025 & 0.0025 \\
3135 & Technicians in control of metal production processes & 0.000 & 0.0025 & 0.0024 \\
7891 & Divers & 0.000 & 0.0024 & 0.0025 \\
2121 & Non-specialized nurses & 0.000 & 0.0023 & 0.0025 \\
3716 & Lay auxiliaries of religions & 0.112 & 0.0023 & 0.0025 \\
2123 & Midwives & 0.052 & 0.0025 & 0.0024 \\
2112 & Other medical specialists & 0.000 & 0.0024 & 0.0025 \\
1111 & Members of the executive branch (national, regional and local) and the legislative branch & 0.000 & 0.0025 & 0.0000 \\
1312 & Production directors of fishing and aquaculture farms & 0.000 & 0.0024 & 0.0000 \\
5992 & Bathers-lifeguards & 0.000 & 0.0021 & 0.0024 \\
6430 & nan & nan & 0.0000 & 0.0024 \\
1314 & Directors of mining operations & 0.013 & 0.0024 & 0.0000 \\
3623 & Professionals supporting public administration licensing services & 0.102 & 0.0023 & 0.0000 \\
4444 & Pawn and loan shop employees & 0.065 & 0.0022 & 0.0000 \\
3632 & Civil Guard non-commissioned officers & 0.037 & 0.0022 & 0.0000 \\
2591 & Notaries and registrars & 0.022 & 0.0021 & 0.0000 \\
3631 & National, regional and local police technicians & 0.037 & 0.0021 & 0.0000 \\
\midrule
\multicolumn{5}{p{0.93\textwidth}}{\footnotesize \emph{Notes:} Panel A lists occupations above the fixed 0.1169 cutoff. Panel B is ordered by the larger of the two outcome-specific weights and omits donors receiving zero weight in both estimations. Occupation descriptions are English translations of official CNO titles. Donor weights sum to one separately by outcome.} \\
\bottomrule
\end{longtable}
\normalsize


\end{document}
